\documentclass[11pt]{article}

\usepackage{amsmath,amssymb,amsthm}
\usepackage[margin=1.05in]{geometry}
\usepackage{bm}
\usepackage{booktabs}
\usepackage{array}
\newcolumntype{P}[1]{>{\raggedright\arraybackslash}p{#1}}
\usepackage[colorlinks=true,linkcolor=blue,citecolor=blue,urlcolor=blue]{hyperref}

\input{marcos.inc}

\numberwithin{equation}{section}
\sloppy
\setlength{\emergencystretch}{3em}
\newtheorem{proposition}{Proposition}[section]
\newtheorem{lemma}[proposition]{Lemma}
\theoremstyle{remark}
\newtheorem{remark}[proposition]{Remark}

% --- notation (follows doc/slip_interface.tex) -----------------------------
\newcommand{\R}{\mathbb{R}}
\newcommand{\eps}{\varepsilon}
\newcommand{\dS}{\,\mathrm{d}S}
\newcommand{\dV}{\,\mathrm{d}V}
\newcommand{\tr}{\operatorname{tr}}
\newcommand{\cof}{\operatorname{cof}}
\newcommand{\Div}{\operatorname{Div}}
\newcommand{\sym}{\operatorname{sym}}
\newcommand{\dev}{\operatorname{dev}}
\newcommand{\DtN}{\mathrm{DtN}}
\newcommand{\vu}{\mathbf{u}}
\newcommand{\vv}{\mathbf{v}}
\newcommand{\vw}{\mathbf{w}}
\newcommand{\vm}{\mathbf{m}}
\newcommand{\vn}{\hat{\mathbf{n}}}
\newcommand{\vN}{\mathbf{N}}
\newcommand{\vs}{\mathbf{s}}
\newcommand{\vb}{\mathbf{b}}
\newcommand{\vg}{\mathbf{g}}
\newcommand{\vx}{\mathbf{x}}
\newcommand{\vt}{\mathbf{t}}
\newcommand{\bnu}{\bm{\nu}}
\newcommand{\Ch}{\hat{C}}
\newcommand{\Se}{\mathbf{S}_e}
\newcommand{\MS}{M_{\mathrm{S}}}
\newcommand{\MF}{M_{\mathrm{F}}}
\newcommand{\SF}{\Sigma_{\mathrm{F}}}
\newcommand{\code}[1]{\texttt{#1}}

\title{The quasi-static problem classes of \texttt{mfemElasticity}:\\
       physical models, approximations, discretisation and solvers}
\author{Compiled from the repository's method notes and class documentation}
\date{}

\begin{document}
\maketitle

\begin{abstract}
\noindent
This document collects, for each linearised quasi-static problem class of
the library, the physical model and its weak form, every physical
assumption or approximation made beyond linearisation, the finite-element
discretisation, and the numerical methods used to solve the discrete
system. The classes are ordered from the simplest to the richest:
non-gravitating traction and clamped problems (with a general, possibly
pre-stressed and non-natural reference state supplied through the
rheology); the mixed referential--spatial self-gravitating problem in
Dahlen's organisation, including the full ladder of approximate
core--mantle-boundary conditions used in glacial isostatic adjustment
(GIA) codes together with their measured cost; the same problem with the
fluid treated in the gauged (relabelling) formulation; the fully
referential self-gravitating problem about a general equilibrium
(an arbitrary elastic tensor, equilibrium stress and equilibrium mapping
of the reference body); and the problem with a genuinely slipping
fluid--solid interface, in its single-valued and broken-potential
organisations and with both constraint-enforcement paths
(penalty with augmented-Lagrangian iterations, and a monolithic
Karush--Kuhn--Tucker (KKT) saddle-point system). A short closing section describes the viscoelastic time layer
that runs on top of the elastic problems. Methods reused across classes
--- coupling between a mesh and a sub-mesh of it, the
Dirichlet-to-Neumann (exterior-gravity) boundary condition, projected
Krylov solvers and rigid-mode handling, the penalty-and-refinement gauge
fixing, the mapping layer --- are described in detail at first use only.
The content is taken from the repository's notes
(\code{doc/*.md}, \code{doc/*.tex}) and header documentation; points where
the sources disagree or are unclear are listed in a comment block at the
end of the source file.
\end{abstract}

\tableofcontents

%==========================================================================
\section{Common framework}
\label{sec:common}
%==========================================================================

\subsection{Notation and vocabulary}
\label{sec:notation}

This subsection fixes every symbol and term used before its detailed
treatment; later sections define their own local symbols at first use.

\paragraph{Geometry.} $d\in\{2,3\}$ is the space dimension. $M$ is the
(reference) body, $\MS$ and $\MF$ its solid and fluid regions (in
\S\ref{sec:slip} written $M_{\mathrm{s}}$, $M_{\mathrm{f}}$), and
$\Sigma$ the fluid--solid interfaces. $B\supset M$ is the computational
ball of radius $b$, whose outer sphere $\partial B$ (the \emph{DtN
sphere}) carries the exterior condition; the shell $B\setminus M$ is the
(vacuum) \emph{buffer}. CMB and ICB are the core--mantle and inner-core
boundaries.

\paragraph{Fields.} Vectors and vector fields are bold. $\vu$ is the
displacement and $\vv$ (or $\vu'$) a test function. $\rho$ is the
density, $G$ the gravitational constant, $\Phi_0$ the background
gravitational (plus centrifugal) potential and $g = |\grad\Phi_0|$.
$\phi$ is the \emph{Eulerian} (spatial) perturbation of the gravitational
potential, used by the mixed class, with test function $\chi$ (or
$\phi'$); $\zeta = \Phi\circ\varphi$ is the \emph{referential} potential
(the spatial potential $\Phi$ composed with the motion), with background
$\zeta^0$ and perturbation $\zeta^1$, used by the referential classes.
$\psi$ is an applied (tidal) potential and $\sigma$ a surface mass load.
The reference state is described by the triple $(\Ch,\Se,\varphi_e)$:
the second elastic tensor $\Ch$, the equilibrium second Piola--Kirchhoff
stress $\Se$ and the equilibrium mapping $\varphi_e$, all defined in
\S\ref{sec:refstate}. A \emph{rheology} is the library object that owns the
material data (elastic tensor, and for viscoelastic bodies the relaxation
branches of \S\ref{sec:visco}) and supplies the stiffness integrators.

\paragraph{Operators and the macros of \code{marcos.inc}.} $\grad$ is the
gradient and $\divg$ the divergence of a vector field; $D\vu$ is the
derivative (Jacobian matrix) of a vector field, $(D\vu)_{iA} = \partial
u_i/\partial x_A$; $\eps(\vu) = \sym D\vu$ is the linearised strain, with
$\sym A = \frac12(A + A^T)$, $\dev A = A - \frac1d(\tr A)\mathbf{1}$, $\tr$
the trace and $\mathbf{1}$ the identity. $\Div$ is the referential
divergence of a tensor field (row-wise), $\cof F = (\det F)F^{-T}$ the
cofactor matrix, $A:B = A_{ij}B_{ij}$ the double contraction,
$\langle\cdot,\cdot\rangle$ an inner product or duality pairing, and
$\circ$ composition. $\jump{\cdot}$ denotes the jump of a quantity across
an interface (solid side minus fluid side), $\norm{\cdot}$ the Euclidean
norm of a coefficient vector (or the induced matrix norm), and
$\abs{\cdot}$ an absolute value, Euclidean length or Frobenius norm.
$\dS$ and $\dV$ are surface and volume measures. Lower-case indices
$i,j,k,l$ are physical, upper-case $A,B$ referential; repeated indices are
summed. A prime on a field ($\vu'$, $\phi'$) denotes a second argument of a
bilinear form, except in $\rho'_F$ (\S\ref{sec:mixed}).

\paragraph{The Mandel convention.} A symmetric $d\times d$ tensor $e$ is
stored as the vector $\hat e\in\R^{n_s}$, $n_s = d(d+1)/2$, of its diagonal
components and its off-diagonal components multiplied by $\sqrt2$, so that
$\hat e\cdot\hat f = e:f$. A fourth-order tensor $C$ with the minor and
major symmetries is stored as the symmetric $n_s\times n_s$ matrix $\hat C$
with $\hat C\hat e\cdot\hat f = (Ce):f$ (21 independent components in 3-D,
6 in 2-D). ``Mandel form'' and a hat over a tensor-valued expression refer
to this convention (a hat over the bold $\mathbf{n}$, $\vn$, is instead a
unit normal).

\paragraph{Discretisation vocabulary (MFEM).} Fields are discretised in
$H^1$-conforming Lagrange finite-element spaces; a \emph{dof} is a
coefficient of the discrete field. For a vector field the \emph{vdofs} are
the scalar coefficients of all components, and an \emph{L-vector} is the
per-process vector of vdofs (the layout of an assembled linear form before
boundary conditions). \emph{True dofs} are the globally unique unknowns
after shared dofs are identified across processes; a prolongation
$P_{\mathrm{T}}$ maps true dofs to L-vectors, and linear solves are posed
on true dofs. Elements and boundary faces carry integer
\emph{attributes} (region labels); a \emph{marker} is a Boolean array
selecting attributes. A \emph{SubMesh} is a mesh built from a subset of
the elements of a \emph{parent} mesh (or from a subset of its boundary
faces), inheriting their attributes; a SubMesh may be disconnected. A
\emph{shadow space} (\S\ref{sec:submesh}) is the parent's finite-element
collection placed on a SubMesh. \emph{Isoparametric} means the geometry is
represented with the same polynomial order as the field.

\paragraph{Solvers.} CG is the conjugate-gradient method (symmetric
positive definite systems), MINRES the minimal-residual method (symmetric,
possibly indefinite systems, with a symmetric positive definite
preconditioner); both are Krylov methods. GS is a Gauss--Seidel sweep and
AMG algebraic multigrid (BoomerAMG from \emph{hypre}). A preconditioner, an
approximate inverse, is written $\mathcal{M}$. The Lanczos process is used
to estimate extreme eigenvalues.

\paragraph{Abbreviations.} GIA: glacial isostatic adjustment. PREM: the
Preliminary Reference Earth Model. DtN: Dirichlet-to-Neumann
(\S\ref{sec:dtn}). AL: augmented Lagrangian. KKT: Karush--Kuhn--Tucker
(saddle-point) system. Frequently cited papers are abbreviated WD
(Woodhouse \& Deuss \cite{WoodhouseDeuss2007}), AC18 (Al-Attar, Crawford,
Valentine \& Trampert \cite{AC18}), AW10 (Al-Attar \& Woodhouse
\cite{AW10}) and MA24 (Maitra \& Al-Attar \cite{MaitraAlAttar2024}).
``Dahlen's treatment'' of a fluid region means the elimination of the
fluid displacement in favour of the potential (\S\ref{sec:mixed});
``gauged'' means the alternative in which the fluid keeps its displacement
and a gauge is fixed numerically (\S\ref{sec:gauged}).

\paragraph{Overloaded symbols.} A few letters are used in more than one
sense, always distinguished by context and defined locally: $\lambda$
(Lam\'e parameter; generalised eigenvalue; interface multiplier), $P$
(first Piola--Kirchhoff stress; projector; constraint-penalty matrix;
dof injection), $M$ (body; mass matrices), $h$ (element size; Love number;
$dt/\tau$), $Q$ (gauge-penalty matrix), $\eps$ ($\eps(\vu)$ with an
argument is the strain, a bare $\eps$ a small parameter), $\sigma$ (Cauchy
stress; surface mass load; slip map in \S\ref{sec:slip}), $S$ (second
Piola--Kirchhoff stress; Schur complements), $B$ (the ball; various
matrices $B_u$, $B_n$, $B_q$), $C$ (elastic tensors; the coupling matrix of
\S\ref{sec:mixed}), $K$ (Laplacian stiffness matrix; a number of
iterations; calligraphic $\mathcal{K}$ for the relabelling kernel and the
block operator), $\kappa$ (bulk modulus; condition number $\kappa(\cdot)$),
and $w$
(non-bold: an augmented-Lagrangian multiplier vector; bold $\vw$: a
displacement-like field defined where used).

\subsection{Scope}

All classes solve the \emph{linearised, quasi-static} equations of a
(possibly self-gravitating, possibly pre-stressed) elastic body: the
inertial term $\partial_t^2\vu$ and the Coriolis term are dropped, and
rotation enters only through the centrifugal background, which for a
constant rotation vector $\bm{\Omega}$ is absorbed into the background
potential, $\Phi_0 \to \Phi_0 + \psi_{\mathrm{c}}$ with
$\psi_{\mathrm{c}}$ the centrifugal potential (AC18; MA24). Time enters only through the data (loads,
and, in the viscoelastic layer of \S\ref{sec:visco}, internal variables).
Every class implements one small interface: at a time $t$,
\code{AssembleForce}$(t)$ brings all time-dependent data to $t$ and
clears increments, \code{AddForce} superposes dual vectors on the
displacement row (in the L-vector layout), and \code{Solve} returns the displacement (and any
further unknowns, which remain internal). There is one displacement field
per class (two in the slip class); several solid regions share it on a
possibly disconnected SubMesh, regional material differences being the
rheology's business. Serial and parallel execution are handled by the same
class: the finite-element space decides.

\subsection{Naming and the problem-class grid}

Class names are built from fixed axis slots,
\begin{multline*}
  [\text{Linear}|\text{Nonlinear}]\;[\text{QuasiStatic}|\text{Dynamic}]\;
  [\text{Mixed}|\text{Referential}]^{?}\\
  [\text{SelfGravitating}]^{?}\;[\text{variant}]\;\text{Problem}.
\end{multline*}
The elasticity is referential in \emph{every} class; the formulation slot
names how the \emph{gravity} is described, and so exists only when the
problem self-gravitates. \emph{Mixed} is the referential--spatial
organisation (referential displacement, Eulerian potential perturbation:
Dahlen's arrangement); \emph{Referential} is the fully referential one.
Properties of the \emph{reference state} are carried by the rheology, not
by the names (\S\ref{sec:refstate}). Table~\ref{tab:grid} summarises the
classes treated here.

\begin{table}[h]
\centering
\footnotesize
\begin{tabular}{@{}P{4.4cm}P{2.0cm}P{2.8cm}P{2.6cm}P{2.6cm}@{}}
\toprule
Class (prefix \code{LinearQuasiStatic}) & Gravity & Reference state & Fluid regions & Solver \\
\midrule
\code{TractionProblem}, \code{ClampedProblem} (\S\ref{sec:nograv})
  & none & general $(\Ch,\Se,\varphi_e)$ via rheology & optional gauged & (projected) CG \\
\code{MixedSelfGravitating\-Problem} (\S\ref{sec:mixed}, \S\ref{sec:gauged})
  & Eulerian $\phi$ & hydrostatic and natural, \emph{required} & Dahlen (\code{FluidRegion}), or gauged & block MINRES; Schur CG \\
\code{ReferentialSelf\-GravitatingProblem} (\S\ref{sec:referential})
  & referential $\zeta^1$ & general & welded gauged & block MINRES \\
\code{ReferentialSelf\-GravitatingSlipProblem} (\S\ref{sec:slip})
  & referential, single-valued or broken & general (broken organisation for $\varphi_e\neq\mathrm{id}$) & slipping, gauged & block MINRES with AL sweeps; or KKT MINRES \\
\bottomrule
\end{tabular}
\caption{The quasi-static problem classes. ``Schur CG'': CG on the
displacement Schur complement (\S\ref{sec:mixed-solvers}); ``AL sweeps'':
penalty with augmented-Lagrangian iterations (\S\ref{sec:al}); ``welded'':
one continuous displacement across solid and fluid (\S\ref{sec:gauged});
``single-valued or broken'': one potential on the ball, or one per region
(\S\ref{sec:slip}).}
\label{tab:grid}
\end{table}

\subsection{Kinematics, stress measures and reference states}
\label{sec:refstate}

We follow the notation of Maitra \& Al-Attar \cite{MaitraAlAttar2021}
(close to Marsden \& Hughes \cite{MarsdenHughes}). A reference body $M$
with referential coordinates $\vx$ is mapped by a motion $\varphi$ into
physical space; $F = D\varphi$ is the deformation gradient (second index
referential), $J = \det F$ its Jacobian, and $F^TF$ the right
Cauchy--Green tensor. With $\varrho$ the physical (spatial) density, the
referential density is $\rho = J\,\varrho\circ\varphi$. A reference
surface element with unit normal $\vN$ and area $\dS$ is carried to a
physical element with unit normal $\vn$ and area $\mathrm{d}S_{\mathrm{phys}}$
related by \emph{Nanson's relation}
\begin{equation}
  \vn\,\mathrm{d}S_{\mathrm{phys}} = \cof(F)\,\vN\dS = J F^{-T}\vN\dS
  =: \bnu\dS, \qquad \mathrm{d}S_{\mathrm{phys}} = |\bnu|\dS,
  \label{eq:nanson0}
\end{equation}
which defines the unnormalised \emph{Nanson normal} $\bnu$. The material
behaviour is given by a stored-energy density $W(\vx,F) = V(\vx,F^TF)$
(per referential volume); $D_F$ denotes the derivative with respect to
$F$. The first and second Piola--Kirchhoff stresses are $P = D_F W$ and
$S = 2 D_{F^TF} V$ (the Cauchy stress being $\sigma = J^{-1}PF^T$), and
there are two elastic tensors,
\begin{equation}
  A = D_F^2 W \ \ \text{(major symmetry only)}, \qquad
  \Ch = 4 D_{F^TF}^2 V \ \ \text{(classical symmetries, 21 components)} .
\end{equation}
An \emph{equilibrium} $\varphi_e$ (with $F_e = D\varphi_e$,
$J_e = \det F_e$, and $P_e$, $S_e$, $\sigma_e$ its stresses) satisfies
$\Div P_e + \rho\bm{\gamma}_e = 0$, with $\bm{\gamma}_e$ the gravitational
plus centrifugal acceleration; nothing requires the equilibrium stress to
be hydrostatic.

Following Al-Attar \& Crawford \cite{AlAttarCrawford2016}, two
independent properties of the reference state are distinguished:
\begin{itemize}
\item \textbf{hydrostatic} versus non-hydrostatic: the equilibrium second
  Piola--Kirchhoff stress is $\Se = -p^0\mathbf{1}$ with $p^0$ the
  equilibrium pressure, or general;
\item \textbf{natural} versus non-natural: the particle label is the
  equilibrium position, $\varphi_e = \mathrm{id}$, or any other labelling,
  $\varphi_e \neq \mathrm{id}$.
\end{itemize}
At a natural reference $P_e = S_e = \sigma_e =: T^0$ (the initial
stress) and
\begin{equation}
  A_{iAjB} = \Ch_{iAjB} + \delta_{ij}\,T^0_{AB} .
  \label{eq:AfromC}
\end{equation}
The reference body is arbitrary: physical statements are invariant under
\emph{particle relabellings}, diffeomorphisms $\xi$ of the reference body
(with $F_\xi = D\xi$, $J_\xi = \det F_\xi$) that change the labels without
changing the physical state, under which (AC18, eqs.~83, 134--136)
\begin{equation}
  \tilde\rho = J_\xi\,\rho\circ\xi, \qquad
  \tilde W(\tilde\vx,\tilde F) = J_\xi\,W(\xi(\tilde\vx), \tilde F F_\xi^{-1}),
  \qquad
  \tilde A_{ijkl} = J_\xi (F_\xi^{-1})_{jm}(F_\xi^{-1})_{ln} A_{imkn}\circ\xi,
  \label{eq:relabel}
\end{equation}
with surface multipliers transforming with the \emph{surface} Jacobian.
In the library a rheology carries the triple $(\Ch, \Se, \varphi_e)$
(\code{ReferentialElasticRheology}); hydrostatic and natural states are
special \emph{values} of it.

\subsection{The quadratic energy and the moduli dictionary}
\label{sec:dictionary}

About a natural reference the quadratic elastic-plus-initial-stress energy
is (WD, eqs.~12--16; AC18 eq.~97)
\begin{equation}
  E_2(\vu) = \frac12\int_M \Ch\,\eps(\vu):\eps(\vu)
           + T^0_{AB}\,\partial_A u_k\,\partial_B u_k \dV ,
  \label{eq:E2}
\end{equation}
the initial stress acting on the \emph{full} gradient. About a general
equilibrium mapping the material term uses the linearised Green strain
(the Green strain is $\frac12(F^TF - \mathbf{1})$; its linearisation about
$F_e$ in the direction $D\vu$ is $\sym(F_e^TD\vu)$), and the implemented
total-Lagrangian split (all quantities referred to the reference body) is
\begin{equation}
  a_{\mathrm{ref}}(\vu,\vv) = \int_M
     \bigl\langle \Ch\,\widehat{\sym(F_e^TD\vu)},\,
                       \widehat{\sym(F_e^TD\vv)} \bigr\rangle
   + \Se : (D\vu^T D\vv) \dV ,
  \label{eq:matgeo}
\end{equation}
(hats: Mandel vectors, \S\ref{sec:notation}), the \emph{material} stiffness carrying no Jacobian
(strain energy per referential volume) and the \emph{geometric} stiffness
$\int S_{AB}\,\partial_A u_k\,\partial_B v_k$ carrying no $F_e$.

\paragraph{The dictionary.} Seismological models do not tabulate $\Ch$.
With hydrostatic $T^0 = -p^0\mathbf{1}$, WD (eqs.~59--64) rearrange the equations so that
the pressure parts of $A$ combine with $\grad p^0 = -\rho\grad\Phi_0$ into
the familiar gravity terms $\rho[\grad(\vu\cdot\grad\Phi_0) -
\grad\Phi_0\,\divg\vu]$, leaving a stiffness with the \emph{effective}
tensor
\begin{equation}
  C^{\mathrm{eff}}_{ijkl} = \Ch_{ijkl}
    + p^0(\delta_{ij}\delta_{kl} - \delta_{il}\delta_{jk} - \delta_{ik}\delta_{jl}).
  \label{eq:Ceff}
\end{equation}
PREM's moduli are the isotropic components of $C^{\mathrm{eff}}$, not of
$\Ch$; with $\kappa$ the bulk modulus, $\mu$ the shear modulus and
$\lambda = \kappa - 2\mu/3$ the Lam\'e parameter, the conversion is
\begin{equation}
  \lambda_{\mathrm{bare}} = \lambda_{\mathrm{PREM}} - p^0, \qquad
  \mu_{\mathrm{bare}} = \mu_{\mathrm{PREM}} + p^0, \qquad
  \kappa_{\mathrm{bare}} = \kappa_{\mathrm{PREM}} - p^0/3 .
  \label{eq:bare}
\end{equation}
This is not small: at the CMB $p^0 \approx 136$~GPa against
$\mu_{\mathrm{PREM}} \approx 290$~GPa. In a fluid (stored energy a
function of $J$ alone, $W = V(\vx,J)$)
$\mu_{\mathrm{PREM}} = 0$ but $\mu_{\mathrm{bare}} = p^0$: the bare tensor
of a fluid at pressure has pressure-sized shear components, which combine
with the explicit $T^0$ term of \eqref{eq:AfromC} to make the incremental
traction purely normal. Consequences used throughout:
\begin{itemize}
\item the \emph{mixed} class (\S\ref{sec:mixed}) assembles $C^{\mathrm{eff}}$
  plus the rearranged gravity terms (WD eq.~59): PREM moduli are correct
  there as they stand;
\item every class assembling \eqref{eq:matgeo} must be given the
  \emph{bare} $\Ch$; the conversion is done at coefficient level
  (\code{BareElasticTensorCoefficient}, taking $p^0$ as a coefficient), so
  model files keep seismological moduli;
\item the two assemblies agree for hydrostatic $T^0$, which is tested as
  an operator identity.
\end{itemize}
The stress dependence of the moduli themselves (the tensor $\Pi$ of
\cite{MaitraAlAttar2021}, describing how the moduli change with
incremental stress) is a constitutive question about Earth models and
does not enter the linearised solvers.

%==========================================================================
\section{Non-gravitating problems: traction and clamped}
\label{sec:nograv}
%==========================================================================

\subsection{Model and weak form}

The classes \code{LinearQuasiStaticTractionProblem} and
\code{LinearQuasiStaticClampedProblem} solve: find $\vu \in V$ with
\begin{equation}
  a(\vu, \vv) = \int_{\Gamma_N} \vt\cdot\vv \dS + \langle f, \vv\rangle
  \qquad \forall \vv \in V_0 ,
  \label{eq:nograv}
\end{equation}
where $V$ is the displacement space ($H^1$ vector fields on $M$, with the
prescribed values on $\Gamma_D$ in the clamped case), $V_0$ the test space
(vanishing on $\Gamma_D$), $\Gamma_N$ the loaded part of the boundary,
$\vt$ the applied traction, $\langle f,\vv\rangle$ any further
(dual-vector) load, and $a$ the stiffness supplied by the rheology (the
object holding the material data):
\begin{itemize}
\item isotropic: $a(\vu,\vv) = \int_M \kappa\,\divg\vu\,\divg\vv +
  2\mu\,\dev\eps(\vu):\dev\eps(\vv)$;
\item anisotropic: $a(\vu,\vv) = \int_M C\eps(\vu):\eps(\vv)$, $C$ in
  Mandel form (21 components in 3-D);
\item referential (\code{ReferentialElasticRheology}): the material +
  geometric split \eqref{eq:matgeo} for a general $(\Ch,\Se,\varphi_e)$ ---
  the non-gravitating referential problem.
\end{itemize}
In the traction problem a traction $\vt$ acts on marked boundary
attributes and there are no essential conditions. In the clamped problem
the displacement is prescribed on $\Gamma_D$ ($\vu = \vu_D$, homogeneous by
default) and a traction is applied on $\Gamma_N$; all other boundaries are
traction-free. (``Mixed'' in the boundary-condition sense only; there is no
relation to the mixed \emph{formulation} of \S\ref{sec:mixed}.) Further
displacement loads $f$ are added through \code{ExternalLoad()} or
\code{AddForce()}.

\subsection{Assumptions and restrictions}

\begin{enumerate}
\item No gravity (no self-gravitation, no background gravity terms) and no
  body force in the operator.
\item The reference state enters only through the rheology. A
  non-hydrostatic or non-natural state is admitted; whether the supplied
  $\Se$ is an equilibrium of the body is the modeller's burden when it is
  supplied as a raw coefficient (the background-state generators of
  \S\ref{sec:background} produce equilibrated fields).
\item The traction is per unit \emph{referential} area. A spatial traction
  on a mapped boundary is composed with the Nanson area factor
  $|\bnu| = \mathrm{d}\tilde S/\dS$ (\code{NansonAreaCoefficient},
  \S\ref{sec:mappings}); composition is the driver's business, coefficients
  stay referential.
\item The pure traction problem requires a load with no net force or
  torque (range compatibility); any incompatible part is projected out
  (\S\ref{sec:projected}).
\item In 2-D the isotropic rheologies describe a two-dimensional continuum
  ($\lambda = \kappa - \mu$), not plane strain of a 3-D body; the
  transversely isotropic tensor coefficient gives the plane-strain
  restriction of the 3-D tensor. The two conventions should not be mixed.
\end{enumerate}

\subsection{Discretisation}
\label{sec:nograv-disc}

The displacement is sought in an $H^1$-conforming vector Lagrange space
on a (possibly curved, isoparametric) mesh. Integrators:
\begin{itemize}
\item the bulk/deviatoric split needs no custom integrator, since
  $\kappa\,\divg\vu\,\divg\vv + 2\mu\,\dev\eps(\vu):\dev\eps(\vv) =
  \lambda\,\divg\vu\,\divg\vv + 2\mu\,\eps(\vu):\eps(\vv)$ with
  $\lambda = \kappa - 2\mu/d$: two MFEM \code{ElasticityIntegrator}s with
  $(\kappa, 1, 0)$ and $(\mu, -2/d, 1)$;
\item anisotropic: \code{ElasticTensorIntegrator}, which at each
  quadrature point forms the Mandel strain--displacement matrix $B_q$
  (mapping an element's vdofs to the Mandel vector $\widehat{\eps(\vu)}$)
  and accumulates $B_q^T\hat C B_q$ times the quadrature weight;
\item referential: \code{MaterialStiffnessIntegrator} (the Mandel vector of
  $\sym(F_e^TD\vu)$ built from the shape gradients and $F_e$; with
  $F_e = \mathbf{1}$ this is \code{ElasticTensorIntegrator}) plus
  \code{GeometricStiffnessIntegrator} (matrix-coefficient vector
  diffusion with $\Se$).
\end{itemize}
The displacement order must be at least the geometry order: a rigid
rotation of a curved element is representable only then. With order-1
displacements on an order-2 mesh the discrete rigid modes are poor null
vectors and results shift at the 10\,\% level on the coarse test problems.

Each (re)assembly builds a fresh bilinear form that borrows the
integrators from a never-assembled template form, rather than calling
\code{Update()} on the old one (which only zeroes the matrix in place, and
whose \code{Finalize(skip\_zeros)} may have dropped entries a new
coefficient needs).

\subsection{Solvers}

After assembly and elimination of essential conditions the discrete
problem is $AX = b$ on true dofs, with $A$ symmetric positive
(semi-)definite. The default solver is preconditioned CG: GS in serial,
BoomerAMG with elasticity options in parallel, warm-started from the
previous solution. Two devices of the base class matter mostly under time
stepping (\S\ref{sec:visco}), where consecutive solves differ little.

\paragraph{Warm starts and the absolute tolerance.} Preconditioned CG
monitors the residual in the preconditioner norm,
$\norm{r}_{\mathcal{M}} := (\mathcal{M}r, r)^{1/2}$, and MFEM stops when
$\norm{r_k}_{\mathcal{M}} \le \max(\tau_{\mathrm{rel}}\norm{r_0}_{\mathcal{M}},
\tau_{\mathrm{abs}})$. From a cold start $X_0 = 0$, $r_0 = b$ and the
target is $\tau_{\mathrm{rel}}\norm{b}_{\mathcal{M}}$. From a warm start
$r_0 = b - AX_0$ may already be $\sim10^{-12}\norm{b}$, so a relative
criterion would ask for $\sim10^{-24}$ and CG would stall at its iteration
limit. The base class therefore makes a warm-started solve converge to the
\emph{same} target as a cold one:
\begin{enumerate}
\item form $\beta = \norm{b}_{\mathcal{M}} = (\mathcal{M}b, b)^{1/2}$ (one
  preconditioner application and one global dot product);
\item if $\beta = 0$, return $X = 0$ without solving;
\item otherwise set $\tau_{\mathrm{abs}} = \tau_{\mathrm{rel}}\beta$ and
  disable the relative criterion, then run CG from $X_0$.
\end{enumerate}
The returned solution then satisfies the cold-start criterion whatever the
quality of the initial guess (the default $\tau_{\mathrm{rel}} =
10^{-12}$).

\paragraph{Preconditioner reuse.} When the operator is reassembled
(relaxation weights, \S\ref{sec:visco}) the preconditioner $\mathcal{M}$,
whose AMG setup dominates the cost, is kept while the iteration count
stays within a factor (default 2) of the count recorded when it was built;
past that it is rebuilt. The matrix is always the current one, so only the
convergence rate, never the solution, depends on the lag.

\subsection{Rigid modes and projected Krylov solvers}
\label{sec:projected}

A free body under traction (and, later, a self-gravitating body) has
rigid motions as (near-)null vectors of its symmetric operator $A$. Two
\emph{different} projections are involved, in different inner products.

\paragraph{The modes.} The rigid motions $\vu_r = \mathbf{a} +
\mathbf{b}\times\vx$ (in 2-D, $\mathbf{b}\times\vx$ is the in-plane
rotation) are interpolated on the displacement space and orthonormalised
in the Euclidean true-dof inner product by modified Gram--Schmidt
(\code{NullSpaceProjector}; a vector numerically dependent on those
already present is dropped, so callers may add modes freely). Collect them
as the columns of $N = [\mathbf{n}_1,\dots,\mathbf{n}_m]$, $N^TN = I$, and
define the Euclidean orthogonal projector
\begin{equation}
  P = I - NN^T, \qquad P = P^T = P^2, \qquad PN = 0 .
  \label{eq:projector}
\end{equation}

\paragraph{Range compatibility.} Since $A = A^T$, $\mathrm{range}(A) =
\ker(A)^\perp$ in the Euclidean inner product. If $N$ spans $\ker A$
exactly, $AX = b$ is solvable iff $N^Tb = 0$ (no net force or torque), and
$Pb$ is the compatible part of any load. For physically consistent loads
replacing $b$ by $Pb$ is a round-off correction.

\paragraph{The projected system.} \code{ProjectedSolver} hands its Krylov
method the operator $PAP$ and solves
\begin{equation}
  PAP\,X = Pb, \qquad X\in\mathrm{range}(P) .
  \label{eq:projsys}
\end{equation}
$PAP$ is symmetric and maps $\mathrm{range}(P)$ into itself, so a Krylov
iteration started in $\mathrm{range}(P)$ stays there; and $PAP$ is
positive definite on $\mathrm{range}(P)$ whenever $A$ is positive on the
complement of the modes. Hence \eqref{eq:projsys} has a unique solution in
$\mathrm{range}(P)$, which CG finds. When $N$ spans $\ker A$ exactly,
$AX = APX = PAPX = Pb$ for $X\in\mathrm{range}(P)$ (because $AN = 0$ and
$A$ symmetric give $A = PAP$), so $X$ solves the original problem. When the
modes are only \emph{near}-null (curved geometry, gravity coupling),
\eqref{eq:projsys} is the problem restricted to the complement of the
modes --- a regularisation, which is the intended meaning in the coupled
classes. The preconditioner is projected in the same way, $P\mathcal{M}P$:
it is symmetric, positive on $\mathrm{range}(P)$, and keeps every
preconditioned residual in $\mathrm{range}(P)$. An unprojected
$\mathcal{M}$ instead re-injects a component along $N$ at every
application; $PAP$ does not see that component, so nothing controls it,
and in practice (BoomerAMG or the shifted Laplacian, in parallel; serial
GS happens not to trigger it) the notes record that this round-off null
component is amplified by a factor ``$1/\eps$'' per iteration (with
$\eps$ the small parameter characterising the near-null direction) and
CG diverges once the residual reaches round-off.

\paragraph{Algorithm (projected solve).}
\begin{enumerate}
\item $b \leftarrow Pb$; $X_0 \leftarrow PX_0$ (a warm start may carry a
  rigid component);
\item run CG (or MINRES) on $PAP$ with preconditioner $P\mathcal{M}P$ from
  $X_0$, with the warm-start tolerance logic above;
\item $X \leftarrow PX$ (removing round-off drift along $N$);
\item optionally re-gauge $X$ (next paragraph).
\end{enumerate}
Iterating on the unprojected $A$ while projecting only data and solution
is weaker on all three counts: the near-null direction remains visible to
the iteration, the preconditioner is unprojected, and the iterate may
acquire a rigid component.

\paragraph{The mass-weighted gauge solve.} Which representative
$X + N\mathbf{c}$ is returned is a separate choice from range
compatibility. The default is the Euclidean one, $N^TX = 0$. The physically
meaningful one is zero net momentum and angular momentum,
$\int\rho\,\vu\cdot(\mathbf{a} + \mathbf{b}\times\vx) = 0$ for all
$\mathbf{a},\mathbf{b}$, i.e.\ $N^T M_\rho X = 0$ with $M_\rho$ the
$\rho$-weighted vector mass matrix (\code{SetMassWeightedGauge}). It is
applied \emph{after} the Euclidean solve: the gauged solution is the unique
$X_g = X + N\mathbf{c}$ with $N^TM_\rho X_g = 0$, i.e.
\begin{equation}
  (N^TM_\rho N)\,\mathbf{c} = -N^TM_\rho X ,
  \label{eq:gaugesolve}
\end{equation}
an $m\times m$ symmetric positive definite system ($m\le6$). Modes of
zero $M_\rho$-norm (the 2-D constant potential of a block system,
\S\ref{sec:mixed-2d}) are excluded from \eqref{eq:gaugesolve} and keep the
Euclidean condition. The inner iteration is deliberately left Euclidean:
an $M_\rho$-orthogonal projector $I - N(N^TM_\rho N)^{-1}N^TM_\rho$ is not
symmetric, and $PAP$ built from it would not be. Strains, and every
rigid-motion-invariant quantity, are the same in both gauges, since $X_g -
X\in\mathrm{range}(N)$.

\paragraph{Non-natural reference states.} With a referential rheology the
rotational modes are those of the \emph{mapped} positions,
$W\varphi_e(\vx)$ with $W$ a constant skew matrix; translations are exact null vectors either
way, the mapped rotations are null through the moment balance of the
background state, and \code{RigidPairResiduals()} measures both
(translations at round-off, rotations decreasing with refinement). The
mass-weighted rotational condition then reads $\int\rho\,\vu\cdot
(\mathbf{b}\times\varphi_e(\vx)) = 0$ with the \emph{referential} density
and no extra Jacobian factor.

The clamped problem needs none of this: the essential conditions remove
the rigid modes.

\paragraph{Fluid regions without gravity.} The base class also accepts
gauged fluid regions (\code{SetGaugedFluid}), used without gravity for
the cavity tests of \S\ref{sec:gauged-verif}; the machinery is described
in \S\ref{sec:gauged}.

%==========================================================================
\section{The mixed self-gravitating problem (Dahlen organisation)}
\label{sec:mixed}
%==========================================================================

The class \code{LinearQuasiStaticMixedSelfGravitatingProblem} couples a
referential displacement on the solid regions to the \emph{spatial}
(Eulerian) gravitational potential perturbation on an enclosing ball, with
a Dirichlet-to-Neumann exterior condition. The weak form follows Al-Attar
\& Tromp \cite{AlAttarTromp2014} (eq.~2.52) and Yu, Al-Attar, Syvret \&
Lloyd \cite{Yu2025} (eq.~3 and Appendix A). Fluid regions are treated in
Dahlen's way: the fluid displacement is eliminated in favour of the
potential. It is the hydrostatic specialisation of the general theory of
\S\ref{sec:referential}.

\subsection{Geometry}

The body $M = \MS \cup \MF$ consists of solid regions $\MS$ (one or more;
e.g.\ inner core and mantle) and fluid regions $\MF$. $B \supset M$ is the
computational ball whose spherical (circular in 2-D) outer boundary carries
the DtN condition. $\Phi_0$ is the background potential; $\sigma$ is a
surface mass load on $\partial M$ (positive when mass is added); $\psi$ is
an applied (tidal) potential. $\SF$ is the union of all fluid--solid
interfaces, $\vm$ \emph{the outward normal of the solid} on it (pointing
into the fluid whether the fluid lies below, as at the CMB, or above, as at
the ICB), and $\rho_F$ the fluid-side density.

\subsection{Assumptions and approximations}
\label{sec:mixed-assume}

\begin{enumerate}
\item[(M1)] \textbf{Hydrostatic reference state}: the equilibrium stress
  is $-p^0\mathbf{1}$, entering only through $\Phi_0$ and the fluid terms
  (the WD rearrangement of \S\ref{sec:dictionary}); the stiffness is
  $C^{\mathrm{eff}}$, so seismological moduli are correct inputs.
\item[(M2)] \textbf{Natural reference state}: $\varphi_e = \mathrm{id}$ ---
  the mesh is the physical equilibrium geometry. (M1) and (M2) are
  requirements of the class, not defaults; the referential class of
  \S\ref{sec:referential} lifts both.
\item[(M3)] \textbf{Fluid regions} are inviscid and barotropic,
  $\rho = \rho(\Phi_0)$ in $\MF$, and every fluid--solid interface is a
  level surface of $\Phi_0$ (both forced by hydrostatic equilibrium).
\item[(M4)] \textbf{Dahlen elimination} of the fluid displacement. In a
  fluid region the linearised balance $-\grad p^1 = \rho^0\grad\phi +
  \rho^1\grad\Phi_0$ (with $\rho^0$ the background density and $p^1$,
  $\rho^1$ the Eulerian pressure and density perturbations) is processed, for perturbations of a \emph{spherical}
  hydrostatic reference, in three moves (the modern rendering is
  AW10, \S\S2.2--2.3): crossing with the radial unit vector
  $\hat{\mathbf{r}}$ shows that $\hat p^1 + \rho^0\phi$ depends on the
  radius $r$ alone, $\hat p^1$ denoting the aspherical part of $p^1$;
  absorbing the degree-0 (spherically symmetric; $l$ denotes the
  spherical-harmonic degree throughout) parts into the spherical reference slaves the aspherical pressure,
  $\hat p^1 = -\rho^0\phi$; a further curl slaves the fluid density,
  $\rho^1 = g^{-1}\partial_r\rho^0\,\phi$ --- the coefficient $\rho'_F$
  below. What remains free is one \emph{constant} pressure perturbation per
  connected fluid region: the degree-0 hole. Consequently the treatment
  never sees the fluid's bulk modulus and misses its compressional physics
  at $l=0$ (\S\ref{sec:cmb-deg0}). About an \emph{aspherical} hydrostatic
  reference the elimination does not readily extend: the free data become
  a whole function $f(\Phi_0)$ per fluid region, coupled to all harmonics
  through Poisson's equation.
\item[(M5)] \textbf{Bare displacement in the gravity coupling.} The
  Eulerian coupling samples the reference potential at displaced positions
  ($\vu\cdot\grad\grad\Phi_0$-type terms). If secular motions are excited
  (true polar wander, slow rotations, drifts that grow displacement without
  strain), this linearisation fails on an aspherical reference; a
  spherically symmetric reference is immune because its potential is
  invariant under the drift \cite{MaitraAlAttar2024}.
\item[(M6)] In 2-D, problems with fluid regions are inconsistent at the
  level of a constant potential (\S\ref{sec:mixed-2d}); 2-D fluid runs are
  for testing only.
\end{enumerate}

\subsection{Weak form}

In the form below $\rho_F$ is the fluid-side density on $\SF$ and
$\rho'_F = \mathrm{d}\rho/\mathrm{d}\Phi_0$ the barotropic density gradient
in the fluid, defined in \eqref{eq:rhoprime}. The symmetric bilinear form
on $(\vu,\phi)\times(\vu',\phi')$ is
\begin{align}
  \mathcal{A}(\vu,\phi\,|\,\vu',\phi') ={}&
    \int_{\MS} \kappa\,(\divg\vu)(\divg\vu') + \int_{\MS} 2\mu\,
      \dev\eps(\vu):\dev\eps(\vu') \nonumber\\
  &+ \tfrac12\int_{\MS}\rho\,\bigl[\grad(\vu\cdot\grad\Phi_0)\cdot\vu'
      + \grad(\vu'\cdot\grad\Phi_0)\cdot\vu\bigr] \nonumber\\
  &- \tfrac12\int_{\MS}\rho\,\bigl[(\vu\cdot\grad\Phi_0)\,\divg\vu'
      + (\vu'\cdot\grad\Phi_0)\,\divg\vu\bigr] \nonumber\\
  &+ \int_{\MS}\rho\,(\grad\phi\cdot\vu' + \grad\phi'\cdot\vu)
   + \frac{1}{4\pi G}\int_{\R^3}\grad\phi\cdot\grad\phi'
  \label{eq:mixedA}\\
  &+ \int_{\MF}\rho'_F\,\phi\,\phi'
  \tag{F1}\\
  &- \int_{\SF}\rho_F\,(\vm\cdot\grad\Phi_0)(\vm\cdot\vu)(\vm\cdot\vu')\dS
  \tag{F2}\\
  &- \int_{\SF}\rho_F\,\bigl[\phi\,(\vm\cdot\vu') + \phi'\,(\vm\cdot\vu)\bigr]\dS .
  \tag{F3}
\end{align}
The first six terms are the solid body; (F1)--(F3) appear with fluid
regions. The elastic terms generalise to an anisotropic tensor without
touching anything else (the gravity, coupling and fluid terms do not
involve the moduli). The $\grad\grad\Phi_0$ term of the Eulerian coupling
has been integrated by parts into the symmetrised form of the second line.

\paragraph{The fluid volume term (F1).} Barotropy gives $\grad\rho =
(\mathrm{d}\rho/\mathrm{d}\Phi_0)\grad\Phi_0$ in the fluid, so
\begin{equation}
  \rho'_F := \frac{\mathrm{d}\rho}{\mathrm{d}\Phi_0}
         = \frac{\grad\rho\cdot\grad\Phi_0}{|\grad\Phi_0|^2}
         \quad\Bigl( = g^{-1}\partial_r\rho \text{ for a radial model}\Bigr).
  \label{eq:rhoprime}
\end{equation}
$\rho'_F < 0$ wherever density increases downward, so (F1) is a
\emph{negative} mass term on the potential (\S\ref{sec:mixed-solvers}).

\paragraph{Interface convention and physical content.} With $\vn$ the
upward unit normal, $\vm\cdot\grad\Phi_0 = -g$ where
the fluid is below and $+g$ where it is above, and (F2)--(F3) reproduce the
published forms, which treat the two kinds of interface separately with
opposite signs. Written with $\vm$ the code never classifies an interface:
every interface integral uses the solid SubMesh's own outward normal, and
the sign of $\vm\cdot\grad\Phi_0$ does the rest. Physically, the
Lagrangian pressure perturbation on the fluid side is $p = -\rho_F(\phi +
\vu\cdot\grad\Phi_0)$ with $\vu\cdot\grad\Phi_0 = (\vm\cdot\vu)
(\vm\cdot\grad\Phi_0)$ on a level surface; the traction $-p\,\vm$ on the
solid gives (F2) and the $\vu'$-half of (F3) (the derivation is
\code{doc/self\_gravitation.md}~\S1). The $\phi'$-half of (F3) is
the surface mass $\rho_F(\vm\cdot\vu)$ displaced across the interface, and
(F1) is the Eulerian density perturbation $\rho'_F\phi$ in the fluid.

\paragraph{Rows and loads.} With the coupling form
$c(\phi,\vu') = \int_{\MS}\rho\,\grad\phi\cdot\vu' - \int_{\SF}\rho_F\,\phi\,
(\vm\cdot\vu')\dS$, $a$ the elastic-plus-gravity displacement form,
$a_\Sigma = $ (F2), $m_F = $ (F1), and the surface-load forms
$\ell_u(\vu') = -\int_{\partial M}\sigma\,\grad\Phi_0\cdot\vu'\dS$,
$\ell_\phi(\phi') = -\int_{\partial M}\sigma\,\phi'\dS$:
\begin{align}
  a(\vu,\vu') + a_\Sigma(\vu,\vu') + c(\phi,\vu') &= \ell_u(\vu') - c(\psi,\vu'),\\
  \frac{1}{4\pi G}\Bigl[\int_B\grad\phi\cdot\grad\phi' + \DtN(\phi,\phi')\Bigr]
   + m_F(\phi,\phi') + c(\phi',\vu) &= \ell_\phi(\phi') - m_F(\psi,\phi') .
\end{align}
After discretisation, with $A_{uu}$ the matrix of $a$, $A_\Sigma$ that of
(F2), $C$ that of $c$ (rows: displacement true dofs, columns: potential
true dofs), $K$ the potential stiffness matrix $\int_B\grad\phi_i\cdot
\grad\phi_j$, $\DtN$ also denoting the DtN matrix (\S\ref{sec:dtn}),
$M_F$ the matrix of $m_F$, and $B_u$, $B_\phi$ the load vectors of
$\ell_u$, $\ell_\phi$, the system is, in blocks,
\begin{equation}
  \begin{pmatrix} A_{uu} + A_\Sigma & C \\ C^T & A_{\phi\phi} \end{pmatrix}
  \begin{pmatrix} \vu \\ \phi \end{pmatrix}
  =
  \begin{pmatrix} B_u - C\Psi \\ B_\phi - M_F\Psi \end{pmatrix},
  \qquad
  A_{\phi\phi} = \frac{K + \DtN}{4\pi G} + M_F ,
  \label{eq:mixedblock}
\end{equation}
with $\Psi$ the interpolant of $\psi$ on the potential space: the tidal
load is the same operators applied to $\Psi$, with no extra assembly.
Without fluid regions the interface and $\rho'_F$ terms are absent and this
is eq.~3 of \cite{Yu2025}. Gravity across interfaces ($[\phi] = 0$,
$[\grad\phi\cdot\vn + 4\pi G\rho\,\vu\cdot\vn] = 0$) is natural in the
weak form \cite{WoodhouseDeuss2007}.

\subsection{Fluid regions and the ladder of CMB approximations}
\label{sec:cmb}

\subsubsection{The condition used in the GIA literature}

Every GIA code in the repository's survey reduces the core to a boundary
condition on the mantle side of the CMB for a \textbf{uniform,
incompressible, inviscid} fluid core, in the lineage Wu \& Peltier
\cite{WuPeltier1982} $\to$ Tromp \& Mitrovica \cite{TrompMitrovica1999}
$\to$ Martinec \cite{Martinec2000} $\to$ Zhong et al.\ \cite{Zhong2003}.
As published:
\begin{itemize}
\item Latychev et al.\ \cite{Latychev2005} (eqs.~11--12), citing Tromp \&
  Mitrovica, in their notation ($T$ stress, $\vs$ displacement, $\Phi^0$
  and $\Phi^1$ background and perturbed potential, superscript $+$ the
  mantle side, $\hat{\bm{\nu}}$ the unit normal out of the mantle into the
  core, $\rho_c$ the core density just below the CMB):
  \begin{equation}
    \hat{\bm{\nu}}\cdot T^+ = \rho_c\,(\vs^+\cdot\grad\Phi^0 + \Phi^1)\,
      \hat{\bm{\nu}}, \qquad
    \sigma_{\mathrm{CMB}} = -\rho_c\,\vs^+\cdot\hat{\bm{\nu}},
  \end{equation}
  traction = buoyancy + potential stress, and an apparent surface mass
  density $\sigma_{\mathrm{CMB}}$ that is the core's \emph{only}
  contribution to $\Phi^1$. The time-marching additionally enforces ``mass
  conservation at the CMB'' each step, standing in for the incompressible
  core's volume constraint.
\item A, Wahr \& Zhong \cite{AWahrZhong2013}: $\sigma_{ij}n_j = -\rho_c\phi\,
  n_i + \rho_c g\,n_i u_jn_j$ (stress $\sigma$, normal $n$, potential
  perturbation $\phi$) --- the same two terms, sign convention differing
  through the normal.
\item Huang et al.\ \cite{Huang2023} (eqs.~2.3.3, 2.2.13):
  $\tau^{\mathrm{CMB}}_{rr} = \rho_c\phi_1 + \rho_c g_0 u_r$ (radial normal
  stress $\tau_{rr}$, potential perturbation $\phi_1$, background gravity
  $g_0$, radial displacement $u_r$; ``potential
  stress'' plus ``Winkler buoyancy''), shear-free; the potential sees the
  core only through the surface density $\rho_c u_r$; ``an incompressible
  and uniform fluid core which has no volumetric density perturbation under
  static deformation''.
\item CitcomSVE-3.0 \cite{CitcomSVE2025}: zero shear force, normal force
  from the self-gravitational effect of a fluid incompressible core
  (Zhong et al.\ 2003); the core is otherwise not considered.
\item The VILMA/Martinec lineage: free-slip CMB on an unmeshed core.
\item The Spada et al.\ \cite{Spada2011} benchmark fixes the Wu \& Ni
  \cite{WuNi1996} solid--fluid conditions as those agreed in the community;
  its reference solutions embed the same physics.
\end{itemize}

\subsubsection{Translation into (F1)--(F3), and the ladder}

With $\vm$ the solid's outward normal, the Latychev/Huang condition is
\emph{exactly} the interface pair (F2) [their $\rho_c g u_r$ buoyancy] and
(F3) [their potential stress and $\sigma_{\mathrm{CMB}}$ source], with two
approximations relative to the full Dahlen treatment: (i) $\rho_F$ is a
\emph{constant} (core-top) rather than the model's fluid-side density on
each interface; (ii) (F1) is absent --- the uniform incompressible core
has no interior density perturbation, $\rho'_F \equiv 0$. These
conditions are degenerate cases of the \code{FluidRegion} machinery,
exposed as a ladder (benchmark option \code{-cmb}):

\begin{center}
\small
\begin{tabular}{@{}lllP{5.6cm}@{}}
\toprule
\code{-cmb} & terms kept & dropped / simplified & \code{FluidRegion} settings \\
\midrule
\code{full} & (F1)+(F2)+(F3) & --- & defaults (stratified fluid) \\
\code{nomass} & (F2)+(F3) & (F1): $\rho'_F = 0$ & \code{density\_gradient} $=0$ \\
\code{uniform} & (F2)+(F3) & (F1); $\rho_F \to$ constant core-top value & $+$ constant \code{interface\_density}: the standard unmeshed-core condition of the GIA literature \\
\code{winkler} & (F2) & (F1); both halves of (F3) & $+$ \code{interface\_potential\_coupling = false}: buoyancy alone (operator stays symmetric), no core-mass contribution to the potential \\
\bottomrule
\end{tabular}
\end{center}

\noindent What is \emph{not} approximate in any rung: the core's
\emph{background} density still generates $\Phi_0$ and enters the
ball-wide Poisson solve. Only the core's \emph{perturbation} physics is
truncated --- as in the published setups, where the core's background
density participates in the zeroth-order gravity. The \code{winkler} rung
is the variant Huang et al.\ call FEMIBF; they find the potential force
``significant for low degree loads''.

\subsubsection{Degree 0 and the core's mass}
\label{sec:cmb-deg0}

For a spherical CMB the surface-density source $\rho_c(\vm\cdot\vu)$ has
zero net mass only when $\oint\vm\cdot\vu\dS = 0$ --- the incompressible
core's volume constraint, living entirely at degree 0 (Latychev et al.\
enforce it per step; other codes leave degree 0 to the surface-load
convention). This is the free constant of (M4) \cite{AW10}: \emph{every}
member of the family, full Dahlen included, misses the core's
compressional physics at $l=0$ because eliminating the fluid discards its
bulk modulus there. The gauged treatment (\S\ref{sec:gauged}) keeps it.

\subsubsection{Measured cost}

\paragraph{Love numbers.} For a spherically symmetric body forced at
spherical-harmonic degree $l$, the surface response is summarised by
dimensionless \emph{Love numbers}: $h$ scales the radial displacement and
$k$ the induced potential perturbation ($l$-Love numbers, scaling the
tangential displacement, are not used here). \emph{Tidal} Love numbers
$h_l, k_l$ refer to forcing by an external potential $\psi$ of degree $l$
(radial displacement $h_l\psi/g$ and induced potential $k_l\psi$ at the
surface); \emph{load} Love numbers $h'_l, k'_l$ refer to forcing by a
surface mass load $\sigma$ of degree $l$, normalised by the potential of
the load itself. Loading responses are concentrated near the surface;
tidal responses sample the deep interior. The external reference is
\code{pyslfp}, an independent radial (one-dimensional) code.

\paragraph{Loading.} Load Love numbers on the \code{prem\_4} model (stratified PREM-like core),
element size $h = 0.2$ (in units of the surface radius), polynomial
order 2; relative change of $h'$ against the full Dahlen
treatment on the same mesh and order, with \code{pyslfp} as the external
reference (absolute values where the comparison to \code{full} is not
meaningful):

\begin{center}
\small
\begin{tabular}{@{}crrrrrr@{}}
\toprule
$l$ & \code{pyslfp} $h'$ & \code{full} & \code{nomass} & \code{uniform} & \code{winkler} & gauged \\
\midrule
0 & $-0.1318$ & $-1.0080$ & $+0.3640$ & $+0.3556$ & $-0.2071$ & $\mathbf{-0.1317}$ \\
1 & $-1.2849$ & $-1.2933$ & 0.08\,\% & 1.4\,\% & 15\,\% & $-1.2839$ \\
2 & $-0.9906$ & $-0.9954$ & 0.08\,\% & 0.08\,\% & 7.9\,\% & $-0.9976$ \\
3 & $-1.0493$ & $-1.0480$ & 0.01\,\% & 0.01\,\% & 2.6\,\% & $-1.0478$ \\
4 & $-1.0511$ & $-1.0470$ & 0.00\,\% & 0.00\,\% & 0.75\,\% & $-1.0468$ \\
\bottomrule
\end{tabular}
\end{center}

\noindent Readings: the standard condition (\code{uniform}) is essentially
harmless at $l\ge2$ for a PREM-like core ($\le0.1\,\%$), and costs
$\sim1.4\,\%$ at $l=1$ (the constant core-top density mis-weights the
response adjacent to the Slichter mode, the gravitationally restored
translation of the inner core, \S\ref{sec:mixed-null}). Dropping the potential stress
(\code{winkler}) is not harmless: 8\,\% at $l=2$, decaying with degree ---
the reason ``potential stress + buoyancy'' became the standard. At $l=0$
the entire family is wrong, each member differently ($h'(0)$ scattered
over $\pm1$ with sign changes), while the gauged treatment agrees with
\code{pyslfp} to $5\times10^{-4}$; no tuning of the CMB condition can
repair a member of the Dahlen family there.

\paragraph{Tides.} A tidal potential samples the deep interior much more
strongly, and the same runs give (relative change against \code{full}):

\begin{center}
\small
\begin{tabular}{@{}cc rrrrrr@{}}
\toprule
$l$ & & \code{pyslfp} & \code{full} & \code{nomass} & \code{uniform} & \code{winkler} & gauged \\
\midrule
2 & $h$ & 0.60412 & 0.60555 & 1.09\,\% & 1.09\,\% & 28.0\,\% & 0.60814 \\
2 & $k$ & 0.29845 & 0.29902 & 2.51\,\% & 2.51\,\% & 42.2\,\% & 0.30043 \\
3 & $h$ & 0.28829 & 0.28693 & 0.21\,\% & 0.21\,\% & 12.4\,\% & 0.28690 \\
3 & $k$ & 0.09215 & 0.09162 & 0.78\,\% & 0.78\,\% & 21.0\,\% & 0.09162 \\
4 & $h$ & 0.17508 & 0.17357 & 0.05\,\% & 0.05\,\% & 5.0\,\% & 0.17351 \\
4 & $k$ & 0.04146 & 0.04102 & 0.26\,\% & 0.26\,\% & 8.7\,\% & 0.04100 \\
\bottomrule
\end{tabular}
\end{center}

\noindent Every approximation costs roughly ten times more on tides than on
loading: \code{nomass}/\code{uniform} reach 1.1\,\% on $h_2$ and 2.5\,\% on
$k_2$, comparable to or above the discretisation error of a production run;
\code{winkler} reaches 28\,\% and 42\,\%. $k$ is consistently more
sensitive than $h$ (about $\times2.3$), being sourced by the internal mass
redistribution that the core treatment controls. \code{nomass} and
\code{uniform} coincide to six digits on tides (for $l\ge2$ the term they
differ by integrates to nearly nothing).

\paragraph{Rotational feedback.} The rotational feedback in sea-level
calculations runs through the degree-2, order-1 components: the centrifugal potential
perturbation is a degree-2, order-1 tidal-type forcing fed back through the
degree-2 \emph{tidal} response. The tidal table is therefore the relevant
error budget. Scaled by the feedback's share of the signal (of order 10\,\%
of the barystatic signal, i.e.\ of the global-mean sea-level change due
to mass exchange between ice and ocean), \code{uniform}/\code{nomass} contribute of order
0.1--0.3\,\% of the total signal (harmless), \code{winkler} of order
3--4\,\% (not harmless). The standard unmeshed-core condition is thus safe
for the loading response and its rotational feedback, the Winkler variant
for neither, and reusing a loading-derived CMB condition for tidal
computations imports a percent-level systematic.

\subsection{Discretisation}

\paragraph{One displacement field on a possibly disconnected SubMesh.} All
solid regions form one \code{SubMesh::CreateFromDomain(parent,} $\{$solid
attributes$\}$\code{)} with one $H^1$ vector space; materials vary by attribute
through the rheology. Nothing in the implementation depends on
connectivity, and BoomerAMG converges on the disconnected space as on a
connected one. The fluid carries no displacement, so the solid regions
never couple to each other directly. The potential $\phi$ is an $H^1$
scalar field on the parent (ball) mesh.

\subsubsection{Coupling forms across a mesh and its SubMesh}
\label{sec:submesh}

The coupling $c(\phi,\vv)$ involves a field on the parent and a field on
the SubMesh, integrated over the SubMesh or its boundary. MFEM's
\code{MixedBilinearForm} indexes trial and test by one element number, and
\code{SubMesh::Transfer} provides no operator a Krylov method could see.
The library's construction (\code{submesh.hpp}) is:

\begin{itemize}
\item \textbf{Principle.} Assemble on the SubMesh with MFEM's own
  \code{MixedBilinearForm}, between the SubMesh space and a \emph{shadow}
  of the parent space (the parent's collection placed on the SubMesh); then
  re-index one side of the sparse matrix from shadow vdofs to parent vdofs.
  The integrand only ever sees one mesh, so every integrator works
  unchanged, including boundary integrators on the SubMesh boundary (the
  interface terms).
\item \textbf{The injection} $\mathcal{I}$ (a matrix of size: number of
  parent vdofs $\times$ number of shadow vdofs, entries $\pm1$, one per
  column; the sign carries orientation bookkeeping) is decoded from MFEM's
  vdof-to-vdof map; sharing the collection \emph{object} guarantees
  identical per-element dof layouts. $\mathcal{I}$ extends by zero,
  $\mathcal{I}^T$ is the exact restriction, $\mathcal{I}^T\mathcal{I} = I$,
  and $\mathcal{I}\mathcal{I}^T$ is the indicator of the parent dofs lying
  in the SubMesh. The re-indexed coupling matrix is $\mathcal{I}C_{\mathrm{sh}}$
  (or $C_{\mathrm{sh}}\mathcal{I}^T$) with $C_{\mathrm{sh}}$ the matrix
  assembled against the shadow space; it is formed as an $O(\mathrm{nnz})$
  re-indexing preserving sparsity, not as a sparse product.
\item \textbf{The true-dof injection $\Pi$ in parallel.} Block operators act
  on true dofs, so the parallel path needs $\Pi$: SubMesh true dofs $\to$
  parent true dofs, as a distributed (hypre) matrix. The tempting
  $\Pi = R_{\mathrm{parent}}\mathcal{I}P_{\mathrm{sub}}$ ($P_{\mathrm{sub}}$
  the SubMesh prolongation, $R_{\mathrm{parent}}$ the restriction of the
  parent L-vector to owned true dofs) is wrong: a shared parent dof on the
  SubMesh boundary may be owned by a process whose local elements there all
  lie outside the SubMesh, and the entry is silently lost (already on a
  $4\times4$ quadrilateral mesh, two slab-partitioned processes, order 2).
  Instead $\Pi^T$ is built row by row over \emph{owned SubMesh} true dofs,
  whose owner always holds a SubMesh element containing them, using the
  global parent true-dof number of the corresponding parent local dof.
  $\Pi$ is Boolean, so $\Pi^T$ is at once the exact primal restriction and
  the dual prolongation.
\item \textbf{Restrictions} (refused by assertion): nonconforming parents,
  variable order, spaces other than $H^1$ and $L^2$, assembly levels other
  than legacy; interior-face integrators. Markers refer to SubMesh
  attributes; cut boundaries receive attribute $\max(\text{parent bdr
  attr})+1$.
\end{itemize}

\subsubsection{The Dirichlet-to-Neumann exterior condition}
\label{sec:dtn}

The \emph{Dirichlet-to-Neumann} (DtN) map of the exterior problem sends
the boundary values of a potential on $\partial B$ to its normal derivative
there, for the harmonic function outside $B$ that has those boundary values
and decays at infinity. Its Galerkin form replaces the unbounded exterior
in the term $\frac{1}{4\pi G}\int_{\R^3}\grad\phi\cdot\grad\phi'$, which
is realised as $\frac{1}{4\pi G}[\int_B\grad\phi\cdot\grad\phi' +
\DtN(\phi,\phi')]$ with $\DtN(\phi,\chi)$ the exterior Dirichlet energy
$\int_{\R^3\setminus B}\grad\tilde\phi\cdot\grad\tilde\chi$ of the decaying
harmonic continuations $\tilde\phi$, $\tilde\chi$ of the boundary traces
(\code{PoissonDtNOperator}).

\paragraph{Harmonic expansion (3-D).} Let $Y_{lm}$ be the real spherical
harmonics of Dahlen \& Tromp \cite{DahlenTromp1998} (Appendix~B),
orthonormal on the unit sphere, with polar angles measured about the
centroid of $\partial B$. A boundary trace expands as $\phi|_{\partial B}
= \sum_{lm}c_{lm}Y_{lm}$ and continues outside as $\tilde\phi =
\sum_{lm}c_{lm}(b/r)^{l+1}Y_{lm}$. Since the outward normal of the exterior
is $-\hat{\mathbf{r}}$ and $-\partial_r\tilde\phi|_{r=b} =
\sum_{lm}\frac{l+1}{b}c_{lm}Y_{lm}$, Green's identity and the orthogonality
$\int_{\partial B}Y_{lm}Y_{l'm'}\dS = b^2\delta_{ll'}\delta_{mm'}$ give
\[
  \int_{\R^3\setminus B}\grad\tilde\phi\cdot\grad\tilde\chi
  = \int_{\partial B}\tilde\chi\,(-\partial_r\tilde\phi)\dS
  = \sum_{lm}(l+1)\,b\,c_{lm}[\phi]\,c_{lm}[\chi].
\]
Writing the coefficients through the boundary moments
$\phi_{lm} := \int_{\partial B}Y_{lm}\phi\dS = b^2c_{lm}$ and truncating at
degree $L$,
\begin{equation}
  \DtN(\phi,\chi) = \frac{1}{b^3}\sum_{l=0}^{L}\sum_{m=-l}^{l}(l+1)\,
     \phi_{lm}\chi_{lm}.
  \label{eq:dtn3}
\end{equation}
In 2-D the expansion is in $\cos k\theta$, $\sin k\theta$, with
arc-length moments $\phi_k = \int_{\partial B}\cos(k\theta)\,\phi\dS$
($k<0$) and $\int_{\partial B}\sin(k\theta)\,\phi\dS$ ($k>0$), $\dS =
b\,\mathrm{d}\theta$, and exterior continuation $(b/r)^{|k|}$. Since
$\int_{\partial B}\cos^2k\theta\dS = \pi b$, the same argument gives
\begin{equation}
  \DtN(\phi,\chi) = \frac{1}{\pi b^2}\sum_{k\neq0}|k|\,\phi_k\chi_k ,
\end{equation}
which is independent of $b$ in terms of the expansion coefficients, as the
2-D exterior Dirichlet energy must be. The $k = 0$ (constant) mode has no
decaying continuation and contributes nothing.

\paragraph{The matrix.} With $\{\phi_i\}$ the potential-space basis
functions that do not vanish on $\partial B$, let $Y$ be the dense
$n_h\times n_\partial$ matrix of boundary moments,
\begin{equation}
  Y_{(lm),i} = \int_{\partial B}Y_{lm}\,\phi_i\dS
  \quad(\text{by boundary quadrature}), \qquad
  \Lambda = \mathrm{diag}\Bigl(\frac{l+1}{b^3}\Bigr),
  \label{eq:dtnmatrix}
\end{equation}
with $n_h = (L+1)^2$ retained harmonics in 3-D ($2L$ in 2-D, where the
moments are the arc-length ones and $\Lambda = \mathrm{diag}(|k|/\pi b^2)$).
Then the DtN
matrix is $Y^T\Lambda Y$ on the boundary dofs, zero elsewhere: symmetric
positive semi-definite, of rank at most $n_h$, coupling every pair of
boundary dofs. Its action costs $O(n_hn_\partial)$: form the $n_h$ moments
$Y x$, scale by $\Lambda$, and scatter back with $Y^T$. In parallel the
moments are partial sums over each process's boundary faces followed by one
global reduction (the operator enters the block system through an RAP
product with the true-dof prolongation). Truncation means that boundary
modes of degree $l > L$ receive no exterior energy (a natural, zero-flux
condition for them); $L$ is the constructor argument \code{dtn\_degree}.
Two consequences recur: (i) the constant has zero DtN energy in 2-D, so the
constant potential is in the kernel of the 2-D Laplace--DtN block (a
potential cannot be pinned at infinity), \S\ref{sec:mixed-2d}; (ii) under
the mappings of \S\ref{sec:mappings} the DtN is untouched, because every
mapping is the identity at and beyond the DtN sphere.

\subsubsection{Assembly of the terms}

\begin{center}
\small
\begin{tabular}{@{}P{2.4cm}P{4.4cm}P{7.4cm}@{}}
\toprule
Term & Lives on & Assembly \\
\midrule
elastic, gravity & solid SubMesh & rheology integrators; symmetrised gravity terms of \eqref{eq:mixedA} \\
(F1) $m_F$ & parent, fluid attributes & \code{MassIntegrator}$(\rho'_F)$ with a domain marker, added to the potential block \\
(F2) $a_\Sigma$ & SubMesh boundary, interface marker & \code{BoundaryNormalNormalIntegrator}$(q)$, $q = -\rho_F(\vm\cdot\grad\Phi_0)$ via \code{BoundaryNormalDotCoefficient}; part of the stiffness integrators \\
(F3) coupling & \code{SubMeshMixed\-BilinearForm} $(\phi\to\vu)$ & \code{BoundaryNormal\-ScalarIntegrator}$(-\rho_F)$ beside the domain term $\rho\grad\phi\cdot\vv$; $C^T$ by transposition \\
tidal load & both rows & $\Psi$ interpolated; $B_u \mathrel{-}= C\Psi$, $B_\phi \mathrel{-}= M_F\Psi$ \\
$\Phi_0$ & parent & Poisson--DtN solve with the solid density (injected from the SubMesh) and fluid densities (on the parent), unless supplied as a coefficient \\
\bottomrule
\end{tabular}
\end{center}

\noindent Details: (a) $\rho'_F$ is supplied analytically
(\code{FluidRegion::density\_gradient}) or evaluated from an element-wise
$L^2$ projection of the density and the discrete $\grad\Phi_0$
(\code{BarotropicDensityGradientCoefficient}); it is assembled after
$\Phi_0$. (b) The solid density is evaluated on the SubMesh; a fluid
density is evaluated on the parent's fluid elements (for $\Phi_0$ and
$\rho'_F$) \emph{and} on the SubMesh's boundary elements on $\SF$ (for
(F2)--(F3)). A coefficient of position serves both; one keyed by domain
attribute does not, since a boundary transformation carries the boundary
attribute --- hence the separate \code{interface\_density}. (c) Interface
integrators take the normal from the boundary transformation, which is
outward from the SubMesh on interfaces inherited from the parent and on cut
ones alike. (d) Under relaxation weights (\S\ref{sec:visco}) only the
displacement block is reassembled (the interface term with it, harmlessly);
potential block, coupling and DtN are built once.

\subsection{Solvers}
\label{sec:mixed-solvers}

Write the system \eqref{eq:mixedblock} as
\begin{equation}
  \mathcal{K}\begin{pmatrix}X_u\\X_\phi\end{pmatrix}
  = \begin{pmatrix}f_u\\f_\phi\end{pmatrix}, \qquad
  \mathcal{K} = \begin{pmatrix} A & C\\ C^T & A_{\phi\phi}\end{pmatrix},
  \qquad A := A_{uu} + A_\Sigma ,
\end{equation}
and let $P$ be the rigid-mode projector \eqref{eq:projector} on the
displacement true dofs ($P_\phi$ its potential counterpart: the identity in
3-D, the projector orthogonal to the constant in 2-D). The block
factorisation
\begin{equation}
  \mathcal{K} = \begin{pmatrix} I & CA_{\phi\phi}^{-1}\\ 0 & I\end{pmatrix}
  \begin{pmatrix} S & 0\\ 0 & A_{\phi\phi}\end{pmatrix}
  \begin{pmatrix} I & 0\\ A_{\phi\phi}^{-1}C^T & I\end{pmatrix},
  \qquad S := A - CA_{\phi\phi}^{-1}C^T ,
  \label{eq:schurfact}
\end{equation}
organises both solvers. $S$ is the displacement Schur complement: the
stiffness of the body once the potential has been eliminated (``the
potential follows the displacement''). For a gravitationally stable body it
is symmetric and positive on the complement of the rigid modes.

\paragraph{Schur CG (reference).} Eliminate the potential and run projected
CG on $S$:
\begin{enumerate}
\item solve $A_{\phi\phi}y = f_\phi$ (inner CG) and form the reduced load
  $g = P(f_u - Cy)$;
\item run projected CG (\S\ref{sec:projected}) on $PSP\,X_u = g$, each
  application of $S$ costing one product with $A$, one with $C^T$, one inner
  CG solve with $A_{\phi\phi}$ (cold-started, relative tolerance a hundredth
  of the outer one and floored at $10^{-13}$, since CG's stopping criterion
  is on the squared residual) and one product with $C$; the preconditioner
  is the projected displacement preconditioner of $A$;
\item recover $X_\phi = A_{\phi\phi}^{-1}(f_\phi - C^TX_u)$ (one more inner
  solve).
\end{enumerate}
The nested inner solves make it roughly an order of magnitude slower than
block MINRES (2-D order 2: 1.1~s against 0.07~s); it is the reference
against which MINRES is checked.

\paragraph{Block MINRES (default).} MINRES on the whole symmetric system,
with the projected block-diagonal preconditioner
\begin{equation}
  \mathcal{M} = \begin{pmatrix} P\mathcal{M}_uP & 0\\ 0 &
     P_\phi\mathcal{M}_\phi P_\phi\end{pmatrix}, \qquad
  \mathcal{M}_u\approx A^{-1}, \quad
  \mathcal{M}_\phi\approx\Bigl(\frac{K + \eps_sM_\phi}{4\pi G}\Bigr)^{-1},
\end{equation}
where $\mathcal{M}_u$ and $\mathcal{M}_\phi$ are GS sweeps (serial) or
BoomerAMG V-cycles (parallel; elasticity options for $\mathcal{M}_u$),
$M_\phi$ is the potential mass matrix and $\eps_s = 10^{-3}$ a shift making
the Laplacian definite. $M_F$ is deliberately left out of
$\mathcal{M}_\phi$ (it may make $A_{\phi\phi}$ nearly singular, below),
and so is the DtN (a preconditioner need only be spectrally equivalent).
MINRES requires only a symmetric operator and a symmetric positive
definite preconditioner, so it applies whatever the inertia of
$\mathcal{K}$; by \eqref{eq:schurfact} that inertia is the inertia of $S$
plus that of $A_{\phi\phi}$. One MINRES iteration costs one product with
each block and one preconditioner application --- no inner solves.

\paragraph{Block CG (diagnostic).} For a gravitationally stable body $S$ is
positive on the complement of the rigid modes and $A_{\phi\phi}$ is
positive (\S\ref{sec:mixed-solvers}, definiteness paragraph), so by
\eqref{eq:schurfact} the projected block system is congruent to a
positive-definite one: the saddle sign of the physical functional does not
survive the symmetrised assembly. CG is therefore admissible on the same
projected block system with the same preconditioner
(\code{SolverType::BlockCG}); its measured cost is identical to MINRES, and
it is kept as a diagnostic (a CG breakdown would signal indefiniteness).

\paragraph{Why the two agree.} The block solver iterates on
$\mathrm{diag}(P,P_\phi)\,\mathcal{K}\,\mathrm{diag}(P,P_\phi)$ with
$X_u\in\mathrm{range}(P)$, i.e.\ it solves
\[
  PAP\,X_u + PC\,X_\phi = Pf_u, \qquad C^TP\,X_u + A_{\phi\phi}X_\phi = f_\phi
\]
(3-D). Eliminating $X_\phi$ from the second equation gives exactly
$P(A - CA_{\phi\phi}^{-1}C^T)P\,X_u = P(f_u - CA_{\phi\phi}^{-1}f_\phi)$,
the projected Schur system of the reference solver. Both therefore return
the same solution to solver tolerance, with no gauge correction
(\S\ref{sec:mixed-null} explains why the projected vector is $(\vu_r, 0)$
and not the coupled pair). Both solvers warm-start across solves.

\paragraph{No segregated iteration.} The coupling is not weak,
$\rho g R/\mu \approx 2.4$ for the Earth, so a block Gauss--Seidel iteration
between $\vu$ and $\phi$ would need under-relaxation with a contraction
factor depending on that ratio; it is not offered.

\paragraph{Definiteness of the potential block.} $A_{\phi\phi} = (K +
\DtN)/4\pi G + M_F$ with $M_F \le 0$. For a uniform fluid ball of radius
$R$ the Laplace--DtN part first becomes singular against the mass term at
$kR = \pi/2$, $k^2 = 4\pi G|\rho'_F|$. For PREM's outer core
$k R_{\mathrm{CMB}} \approx 1.1$: positive for Earth-like models, but not
by a wide margin, and not for much larger or steeper fluid bodies.
\code{PotentialBlockMinEigenvalue()} (Lanczos) reports the margin; a
negative value means the inner CG solves cannot be trusted. The proximity
to the margin also makes (F1) far from negligible: for a PREM-like
three-layer model dropping $\rho'_F$ changes $\norm{\vu}$ and
$\norm{\phi}$ by a factor of about three.

\subsection{Null space and gauge}
\label{sec:mixed-null}

\paragraph{Coupled near-null pairs.} $\vu_r = \mathbf{a} + \mathbf{b}\times
\vx$ on $\MS$, paired with $\phi_r = -\vu_r\cdot\grad\Phi_0$ on all of space,
is an exact null vector of the continuous $\mathcal{A}$. Discretely it is
only near-null: \code{RigidModeResiduals()} gives $\sim10^{-3}$ at order 1
and $\sim10^{-6}$ at order 2, decreasing with refinement. The system is
therefore regularised by projection (\S\ref{sec:projected}).

\paragraph{What is projected.} Both solvers solve the system restricted to
displacements orthogonal to the rigid modes in the Euclidean true-dof inner
product. The MINRES projector removes $(\vu_r, 0)$ (and $(0,1)$ in 2-D),
\emph{not} the coupled pair $(\vu_r,\phi_r)$: the restricted block system
then has exactly the Schur complement the Schur solver uses, so the two
agree to solver tolerance with no gauge correction. Projecting the coupled
pair and fixing the gauge afterwards injects a residual of size
(rigid-mode residual)$\times$(rigid content of the iterate), which the soft
Slichter mode amplifies into a percent-level disagreement. By default the
displacement has no rigid component and the potential is the discrete
solution of its equation for that displacement; with
\code{SetMassWeightedGauge()} the mass-weighted representative is returned
and the potential follows the re-gauged displacement (a rigid shift moves
the body's potential with it). Every projected solve uses the projected
preconditioner $P\mathcal{M}P$.

\paragraph{Enclosed solid regions.} For a spherically symmetric model a
rotation of a solid inner core alone is an exact null vector (no strain,
$\vu\cdot\grad\Phi_0 = 0$, $\vm\cdot\vu = 0$, $\int\rho(\mathbf{b}\times\vx)
\cdot\grad\phi' = 0$); for an aspherical core it is physically near-null
with a restoring torque of the order of the asphericity.
\code{AddRegionRotations()} projects these on request. A
\emph{translation} of the enclosed region is \emph{not} null --- gravity
restores it (the Slichter mode) --- and must never be projected. It is soft
(Rayleigh quotient $x^TAx/x^Tx\sim10^{-4}\norm{A}$ in the test model), so
discretisation error shows first as a spurious inner-core translation;
refinement, not projection, is the remedy.

\subsection{Two dimensions}
\label{sec:mixed-2d}

The constant potential is a null vector of the Laplace--DtN block
(\S\ref{sec:dtn}). Potential loads are made compatible by subtracting a
uniform flux through the outer boundary: with $\mathbf{1}$ the coefficient
vector of the constant function and $\ell_i = \int_{\partial B}\phi_i\dS$
(so that $\mathbf{1}^T\ell = |\partial B|$), $f_\phi\leftarrow f_\phi -
(\mathbf{1}^Tf_\phi/|\partial B|)\,\ell$, after which
$\mathbf{1}^Tf_\phi = 0$; every potential solve runs on $PA_{\phi\phi}P$ with
the constant projected out, and the block solver restricts the potential
likewise. With fluid regions this is a regularisation, not a null-space
fact: the constant is not null once $M_F\neq0$, and the interface coupling
$-\int\rho_F\phi\,(\vm\cdot\vv)$ is not invariant under a constant shift.
The 2-D fluid problem is inconsistent at the level of a constant, the
regularisation changes the answer at the percent level, and 2-D fluid runs
are for cheap testing only. Three dimensions are unaffected.

\subsection{Verification}

Symmetry of the assembled block operator to round-off and decreasing
rigid-mode and region-rotation residuals (a sign error in (F2) or either
half of (F3) would leave them $O(1)$, so this is the sharpest check of the
interface terms); a uniform tidal field $\psi = \mathbf{a}\cdot\vx$, which
loads exactly a rigid mode in the continuum, gives a displacement of the
order of the rigid-mode residual; Schur CG and block MINRES agree to
$10^{-7}$ or better, with and without fluids, serial and parallel. Load
and tidal $h, k$ of a homogeneous sphere against the incompressible
formulas of Wu \& Peltier \cite{WuPeltier1982} with $\kappa/\mu = 200$
agree to 0.1--0.7\,\% at degrees 2--4 on a fine order-2 mesh, the
remainder being the finite bulk modulus. The load's own potential for
$k'$ is taken from a solve for the load alone on the same mesh
(\code{SolveLoadPotential()}: Laplace--DtN without the fluid mass term,
which belongs to the body's response); subtracting the exact value instead
leaves the direct potential's discretisation error in $k'$, a 30\,\% effect
on a coarse 3-D mesh.

%==========================================================================
\section{The mixed problem with a gauged fluid}
\label{sec:gauged}
%==========================================================================

The gauged treatment keeps the fluid displacement in a fully referential
description and fixes the relabelling gauge by a small shear penalty whose
bias is removed by iterated Tikhonov refinement (following Maitra \&
Al-Attar \cite{MaitraAlAttar2024}, \S3.7.3). It is selected in the mixed
class by construction: \emph{no} \code{FluidRegion}s; the fluid attributes
join the displacement SubMesh, the rheology gives them $\kappa_f$ with zero
shear, the density covers them, and \code{SetGaugedFluid()} adds the
penalty. The machinery lives in the base class
(\S\ref{sec:nograv}) and is reused by the referential and slip classes.

\subsection{Model}

\paragraph{The fluid as a solid without shear.} An inviscid fluid region
carries a displacement with stiffness $\kappa_f(\divg\vu)(\divg\vu')$
($\kappa_f$ the fluid's bulk modulus) and zero shear; the gravity, coupling and load terms are the solid's
integrands (the elastic--gravitational form is structure-blind). What
distinguishes a fluid is a symmetry: the solution is defined only up to a
\textbf{linearised relabelling} (a rearrangement of fluid particles with
no Eulerian effect; the displacement is then said to be \emph{gauge},
i.e.\ determined only up to such fields). For a barotropic fluid these are
(Friedman \& Schutz's ``trivial displacements'' \cite{FriedmanSchutz1978})
\begin{equation}
  \mathcal{K} = \{\vw \text{ on } \MF : \divg(\rho\vw) = 0,\
               \vw\cdot\vm = 0 \text{ on the interfaces}\}.
\end{equation}
$\mathcal{K}$ lies in the kernel of the continuous operator, so the fluid
displacement is gauge. The observables are the solid displacement, the
potential, and in the fluid anything built from $\divg(\rho\vu)$ and the
interface normal displacement (e.g.\ the pressure perturbation).

\paragraph{No displacement discontinuity is needed (linear barotropic
problem).} The tangential trace of the fluid displacement is itself gauge,
so using one continuous $H^1$ space over solid and fluid is an admissible
\emph{partial gauge fixing}, not a physical constraint. Given a fluid
solution with tangential slip $\mathbf{j}$ against the solid, extend
$\mathbf{j}$ into the fluid as $\tilde\vw$ with $\tilde\vw\cdot\vm = 0$ and
correct it by $\mathbf{z}$ solving $\divg(\rho\mathbf{z}) =
-\divg(\rho\tilde\vw)$ with vanishing traces (compatible since
$\oint\rho\tilde\vw\cdot\vm = 0$): a relabelling with any prescribed
tangential trace exists. No sphericity is used, so the argument survives
aspherical configurations. It fails for a two-parameter fluid (one whose state depends on a second,
independently advected label such as entropy or composition $s$, which
adds the constraint $\vw\cdot\grad s = 0$) and for the non-linear slip map; those need the slipping interface of
\S\ref{sec:slip}. Two consequences: the tangential traction transmitted is
the fluid's deviatoric stress, $O(\eps)$, removed by the refinement; and an
inner-core rotation now drags the fluid at $O(\eps)$ cost, so
\code{AddRegionRotations()} must \emph{not} be used in this mode (the
penalty owns those modes).

\paragraph{Equivalence with Dahlen and the Adams--Williamson condition.}
The two treatments describe the same physics iff the fluid is
\emph{materially} barotropic (one thermodynamic parameter, so that the
constitutive $\mathrm{d}\rho/\mathrm{d}p$ equals the background's). With
$N^2$ the squared buoyancy (Brunt--V\"ais\"al\"a) frequency of the fluid,
this is
\begin{equation}
  N^2 = 0 \iff \kappa_f = \frac{\rho^2|\grad\Phi_0|}{|\mathrm{d}\rho/\mathrm{d}r|}
  \qquad\text{(Adams--Williamson)} .
\end{equation}
Then the static $\kappa$-response coincides with Dahlen's secular
response. For $N^2\neq0$ the static response is the ``frozen'' elastic one
and differs from Dahlen's by $O(N^2)$, and the gauge freedom shrinks
(\S\ref{sec:slip}), though the penalty still regularises the nearly
neutral case. Benchmarks at $l\ge1$ against Dahlen or \code{pyslfp} must
use neutrally stratified cores (PREM's outer core is Adams--Williamson by
construction; a uniform-density core is not).

\paragraph{Degree 0.} The treatments differ at degree 0 by design: Dahlen's
fluid never sees $\kappa_f$, while the gauged fluid carries it at every
degree, restoring the compressible degree-0 physics (the $l=0$ column of
\S\ref{sec:cmb}).

\paragraph{What changes against Dahlen.}
\begin{center}
\small
\begin{tabular}{@{}P{7.0cm}P{7.0cm}@{}}
\toprule
Dahlen (\S\ref{sec:mixed}) & Gauged \\
\midrule
displacement on the solid SubMesh & displacement on the whole body \\
(F1): $A_{\phi\phi} = (K+\DtN)/4\pi G + M_F$, $M_F\le0$, near-indefinite for steep cores & no (F1): plain Laplace--DtN, SPD with margin \\
interface terms (F2), (F3) & none: interface conditions are natural \\
$c(\phi,\vv) = \int_{\MS}\rho\grad\phi\cdot\vv - \int_\Sigma\rho_F\phi(\vm\cdot\vv)$ & $c(\phi,\vv) = \int_M\rho\grad\phi\cdot\vv$ \\
tidal load on the potential row, $-M_F\Psi$ & none (density perturbation is $-\divg(\rho\vu)$) \\
2-D constant-mode inconsistency & consistent \\
needs $\rho'_F$, level surfaces, interface conventions & needs $\kappa_f$ and the penalty \\
well-conditioned displacement block & $1/\eps$ fluid near-kernel in the displacement block \\
\bottomrule
\end{tabular}
\end{center}

\subsection{Assumptions}

Those of \S\ref{sec:mixed-assume} except (M4) and (M6): the fluid is not
eliminated; instead (G1) the fluid is inviscid, barotropic in the
background and \emph{materially} barotropic (Adams--Williamson) when
equivalence with Dahlen is wanted; (G2) tangential continuity at the
interface is a gauge choice, exact only for the linear barotropic problem;
(G3) no essential boundary conditions on gauged attributes (the penalty's
column elimination is not folded into the loads).

\subsection{Discretisation}

The fluid attributes join the displacement SubMesh. The penalty is one
deviatoric \code{ElasticityIntegrator}$(\eps\mu_g, -2/d, 1)$ on the marked
attributes, i.e.\
\begin{equation}
  q(\vu,\vv) = \int_{\MF} 2\mu_g\,\dev\eps(\vu):\dev\eps(\vv),
\end{equation}
with $\mu_g$ a gauge shear scale (a coefficient) and $\eps$ a small
dimensionless penalty factor, i.e.\ a fluid with shear modulus $\eps\mu_g$;
$q$ vanishes identically on the physical (volumetric) fluid response.
With $Q$ the matrix of $q$ on true dofs, assembly produces, beside the physical $A$ (which \code{SystemMatrix()}
still returns), the penalty $\eps Q$ on true dofs and the regularised
$A + \eps Q$ in one extra assembly (a form borrowing the physical
integrators with the penalty appended). The solver and preconditioner are
built on $A+\eps Q$. $\mu_g$ sets the scale; the fluid's own bulk modulus
is a natural choice. Under a mapping (\S\ref{sec:mappings}) the penalty is
assembled \emph{covariantly}, as the pulled-back deviatoric tensor through
\code{ElasticTensorIntegrator}'s mapped form (\S\ref{sec:ref-gauged}).

\subsection{Gauge fixing by penalty and iterated Tikhonov refinement}
\label{sec:tikhonov}

After discretisation and rigid-mode projection the problem is $Ax = b$,
$A = A^T\succeq0$ ($n\times n$), $b$ compatible (orthogonal to $\ker A$),
with a large \emph{gauge subspace} $\mathcal{G}\subseteq\ker A$ (or nearly
so) of dimension comparable to $n$, not aligned with any set of nodal dofs
(so it cannot be removed by constraining dofs). Discretely the
finite-element space does not contain $\mathcal{K}$ exactly: there is a
large near-kernel of ``almost relabellings'' with eigenvalues of order a
positive power of the element size $h$, and a Krylov method on the raw
system amplifies their small load components into mesh-dependent fluid
displacements polluting the solid through the coupling.

\paragraph{Penalty: what one solve gives.} With $Q = Q^T\succeq0$ positive
on $\mathcal{G}$ (assume for the analysis $Q\succ0$; on the complement of
$\mathcal{G}$ any positive completion leaves the conclusions unchanged),
solve $A_\eps x_\eps = b$, $A_\eps := A + \eps Q$. The generalised eigenpairs
\begin{equation}
  Av_i = \lambda_iQv_i, \qquad v_i^TQv_j = \delta_{ij}, \qquad
  0\le\lambda_1\le\lambda_2\le\dots
  \label{eq:gevp}
\end{equation}
diagonalise $A$ and $Q$ simultaneously ($V^TAV = \Lambda$, $V^TQV = I$).
Gauge directions are the $\lambda_i = 0$ modes; the physical spectrum is
$\{\lambda_i > 0\}$ with smallest member $\lambda_{\min}$. Expanding
$x = \sum_iy_iv_i$ and $b_i := v_i^Tb$ (so $b = \sum_ib_iQv_i$, and
$b_i = 0$ on gauge modes by compatibility), $A_\eps x_\eps = b$ reads
$\sum_i(\lambda_i + \eps)y_{\eps,i}Qv_i = \sum_ib_iQv_i$, while $Ax = b$
reads $\lambda_iy_i = b_i$:
\begin{equation}
  y_{\eps,i} = \frac{b_i}{\lambda_i + \eps}, \qquad
  y_i = \frac{b_i}{\lambda_i}\ (\lambda_i > 0), \qquad
  \frac{y_i - y_{\eps,i}}{y_i} = \frac{\eps}{\lambda_i + \eps}
  \le\frac{\eps}{\lambda_{\min} + \eps} .
\end{equation}
So one penalised solve fixes the gauge exactly (zero component on every
gauge mode: the representative minimising $x^TQx$) but biases each physical
mode by $\eps/(\lambda_i + \eps)$. Shrinking $\eps$ shrinks the bias but
sends the condition number $\kappa(A_\eps)\sim\lambda^A_{\max}/
(\eps\lambda^Q_{\min})$ to infinity, degrading every preconditioner.

\paragraph{Iterated refinement: the recurrence.} Keep $\eps$ moderate and
iterate with the \emph{physical} residual:
\begin{equation}
  x^{(1)} = A_\eps^{-1}b, \qquad
  r^{(k)} = b - Ax^{(k)}, \qquad
  \delta^{(k)} = A_\eps^{-1}r^{(k)}, \qquad
  x^{(k+1)} = x^{(k)} + \delta^{(k)} .
  \label{eq:tikhonov}
\end{equation}
Two identities make it cheap and transparent.
(i) \emph{The residual is the penalty applied to the last increment}:
$r^{(k+1)} = r^{(k)} - A\delta^{(k)} = A_\eps\delta^{(k)} - A\delta^{(k)} =
\eps Q\delta^{(k)}$ (and $r^{(1)} = \eps Qx^{(1)}$). After an exact solve no
product with $A$ is needed: each refinement solves the regularised system
against $\eps Q$ applied to the previous increment.
(ii) \emph{The error contracts geometrically}: with $e^{(k)} = x^{(k)} - x$,
\begin{equation}
  e^{(k+1)} = e^{(k)} - A_\eps^{-1}Ae^{(k)}
            = A_\eps^{-1}(A_\eps - A)e^{(k)}
            = \eps A_\eps^{-1}Q\,e^{(k)},
\end{equation}
which in the basis \eqref{eq:gevp} is diagonal with factor
$\eps/(\lambda_i + \eps)$. Starting from $e^{(1)}_i = -y_i\eps/(\lambda_i +
\eps)$,
\begin{equation}
  \frac{y_i - y_i^{(k)}}{y_i} = \Bigl(\frac{\eps}{\lambda_i+\eps}\Bigr)^k,
  \qquad
  \rho_{\mathrm{contr}} = \frac{\eps}{\lambda_{\min}+\eps}\approx\frac{\eps}{\lambda_{\min}}
\end{equation}
after $k$ solves: the bias $O(\eps)$ of one solve becomes $O(\eps^k)$.
Gauge modes ($\lambda_i = 0$, $b_i = 0$) have factor 1 and stay at their
$Q$-minimal value zero: the iteration never un-fixes the gauge. The fixed
point solves $Ax = b$ for every $\eps > 0$. This is stationary iterated
Tikhonov regularisation \cite{Lardy1975,EnglHankeNeubauer} ($k$-times
iterated Tikhonov has qualification $k$), equivalently defect correction
with a deliberately perturbed operator. If the gauge were a constraint
$Bx = 0$ with $Q = B^TWB$, \eqref{eq:tikhonov} would be the Hestenes--Powell
multiplier iteration with fixed penalty
\cite{Hestenes1969,Powell1969,FortinGlowinski} (\S\ref{sec:al}); the gauge
use is the mirror image (no constraint is wanted; the iteration removes the
penalty's contamination of the physical modes).

\paragraph{Algorithm (gauge refinement, \code{Solve} with a gauged fluid).}
\begin{enumerate}
\item Assemble $A$ and $\eps Q$; build the solver and preconditioner on
  $A_\eps = A + \eps Q$ (\code{RegularizedMatrix()}); the physical $A$ is
  kept for diagnostics.
\item Solve $A_\eps x^{(1)} = b$; set $\delta \leftarrow x^{(1)}$.
\item For $k = 1,\dots,K$ (\code{refinements}, default 2): form
  $r = \eps Q\delta$ and record $\norm{r}$ (\code{GaugeResiduals()}); solve
  $A_\eps\delta = r$ with the same solver and preconditioner; set
  $x\leftarrow x + \delta$.
\end{enumerate}
The ratio of consecutive recorded residuals is the observed contraction
factor. The operator, its preconditioner and Krylov setup are built once,
so refinements cost only additional Krylov solves.

\paragraph{The spectral condition.} Everything rests on $\lambda_{\min}$
being bounded away from zero \emph{uniformly in the discretisation}: on
every physical mode the penalty must be dominated by the physical operator
with a mesh- and order-independent margin. For the gauged fluid this
holds: $Q$ is a deviatoric elasticity form and $A$ contains genuine elastic
stiffness of the same differential order, both scaling identically under
refinement, so $\lambda_{\min}\sim\mu_{\mathrm{solid}}/\mu_g$ and the
observed contraction is $\approx\eps\mu_g/\mu_{\mathrm{solid}}$ per step,
uniformly across meshes and orders. (The condition fails for the ball-wide
vacuum extension of \S\ref{sec:vacuum}.)

\paragraph{Semi-convergence and operating point.} Each application of
$A_\eps^{-1}$ is a Krylov solve to tolerance $\tau$; a residual
perturbation $\eta$ enters the next iterate as $A_\eps^{-1}\eta$, whose
component on mode $i$ is $\eta_i/(\lambda_i + \eps)$ --- amplified by up
to $1/\eps$ on gauge and near-gauge modes. The bias decays like $\rho^k$ only down to a
noise floor of order $\tau/\eps$ (more precisely each step adds
$O(h\text{-residual}/(\eps\mu_g))$ of near-kernel junk); the error as a
function of $k$ is U-shaped. The operating point is
\begin{equation}
  \eps\sim10^{-2}, \qquad k = 2\text{--}3
  \text{ refinements}, \qquad \tau\lesssim10^{-10} ,
\end{equation}
and it is certified from both sides. Above it the $O(\eps)$ bias survives
the refinement budget in the observables (measured: $2.6\times10^{-2}$ in
tidal $h_2$ at $\eps = 10^{-1}$). Below it the cost grows by a measured
$\sim4.5\times$ per decade of $\eps$, and by $\eps = 10^{-4}$ the
iteration is semi-convergent: the gauge component escapes the
refinements and pollutes the solution despite nominally converged solves
($h_2 = -3.17$). The window is problem- and resolution-dependent in
principle; it is demonstrated at PREM structure (\code{prem\_4}, $h = 0.2$,
where the gauged solution matches \code{pyslfp} to $5\times10^{-4}$,
\S\ref{sec:cmb}).

\paragraph{Contraction-rate tripwire.} Since the refinement corrections
contract at $\approx\eps/(\lambda + \eps)$, the ratio of the last two
recorded residuals $\norm{\eps Q\delta}$ diagnoses the operating point.
The mixed class's \code{GaugeRefine()} warns when the ratio exceeds $0.9$
(the semi-convergence signature: $\eps$ too small for this mesh and model)
and prints a note when it exceeds $0.2$ (a residual bias $\sim
\text{rate}^k$ may remain: $\eps$ too large for the refinement budget, or
too few refinements). It only reports; it changes nothing.
Modes adjacent to the projected null space (near-rigid, Slichter-soft
degree-1 responses) have small $\lambda_i$ of physical origin, converge
slowly and dominate the residual diagnostics. The worst case is a load
aligned with a near-null mode: a uniform tidal gradient leaks $\sim3\,\%$ of
a comparable physical response into the solid on the coarse test mesh at
$\eps = 10^{-2}$, $k=3$, growing like $k/\eps$; physical (degree-2) loads do
not sit on the near-kernel.

\paragraph{The coupled (block) version.} In the mixed class the recurrence
runs on the block operator $\mathcal{K}$ of \S\ref{sec:mixed-solvers} with
the penalty $\mathrm{diag}(\eps Q, 0)$ acting on the displacement only, so
identity (i) becomes $r^{(k+1)} = [\eps Q\delta_u^{(k)};\,0]$
(\code{GaugeRefine()} override): each refinement solves the regularised
block system (block MINRES or Schur CG) with zero potential load and no
tidal term, the potential increment is accumulated alongside the
displacement increment (so the potential converges with the displacement
rather than being recomputed), and the accumulated pair is left as the warm
start of the next \code{Solve}.

\subsection{Solvers and null space}

The block solvers of \S\ref{sec:mixed-solvers} run unchanged on $A+\eps Q$,
now with a plain Laplace--DtN potential block. The global rigid modes are
projected as before; inner-core rotations are not projected. The
mass-weighted rigid gauge matters for comparisons against exact radial
solutions: the Euclidean true-dof gauge leaves a mesh-asymmetry rigid
translation ($\sim10^{-5}$ on the test meshes) that uniform refinement does
not remove, and which dominates e.g.\ a Lam\'e comparison at the 0.5\,\%
level.

\subsection{The exact-gauge KKT alternative (parked)}
\label{sec:gaugekkt}

The KKT idea of \S\ref{sec:kkt} applies in principle to the gauge as well:
minimise $\frac12(Qu,u)$ subject to the full coupled equations $Au = f$,
with $Q$ the \emph{unit}-scale deviatoric gauge form on the fluid
attributes. There is then no $\eps$ anywhere: the physical equations hold
exactly, and the gauge is fixed as the minimum-$Q$-energy representative.
With an auxiliary (multiplier) field $v$, determined only up to $\ker A$,
which MINRES on the consistent system tolerates, the optimality system is
the doubled saddle
\begin{equation}
  \begin{pmatrix} \mathrm{diag}(Q, 0) & A\\ A & 0\end{pmatrix}
  \begin{pmatrix} u\\ v\end{pmatrix} =
  \begin{pmatrix} 0\\ f\end{pmatrix},
\end{equation}
with the clean physical $A$ (block operator of \S\ref{sec:mixed-solvers})
in the constraint rows. It is available as \code{EnableGaugeKKT()} on the
mixed class (block organisations only; internally the $A + Q$ assembly of
\code{SetGaugedFluid()} at $\eps = 1$ preconditions both displacement
slots), with \code{KKTResiduals()} reporting the physical residual
$\norm{Au - f}/\norm{f}$ and the stationarity residual $\norm{Qu + Av}$
relative to $\norm{Qu}$.

It is implemented and physically correct, and \emph{parked for cost}. On
the 2-D three-layer gauged test problem and the 3-D fluid-core model at
$h = 0.3$, the iterates reach mesh-level physics early and satisfy
stationarity to $10^{-5}$, but MINRES never converges algebraically: 30k
iterations against the penalty path's $\sim800$ per solve, under
block-diagonal preconditioning at every grade tried --- serial GS, parallel
AMG, and the Murphy--Golub--Wathen composite (on the multiplier slot),
which degrades the solve with smoother-grade components and still caps
with AMG-grade ones. The structural reason: the ``constraint'' operator
$A$ is square and only approximately singular, so saddle-point
preconditioning theory, which assumes a full-rank rectangular constraint,
does not apply, and the near-kernel reappears wherever the scheme needs
solves with $A$. Constraint preconditioning (Keller--Gould--Wathen; MFEM
provides \code{FGMRESSolver} for the required flexible outer iteration) is
the named next rung, but its two inner $A$-solves per outer iteration are
solves of exactly the near-singular system, so even success is expected
to land at penalty-path cost.

The lesson, as \code{doc/kkt\_slip\_solver.md} \S7 closes: KKT pays for a
full-rank interface condition --- the slip constraint, whose multiplier
system reached cost parity with AL (\S\ref{sec:kkt}); for selection over a
near-kernel the penalty is not a compromise but the right
regularisation. The production gauge path therefore remains the
$\eps$-penalty with iterated Tikhonov refinement of \S\ref{sec:tikhonov}.

\subsection{Verification}
\label{sec:gauged-verif}

Gravity-free, exact references: a static fluid supports a uniform pressure,
so its response condenses onto the solid as the rank-one cavity term
$(\kappa_f/V)\,b(\vu)b(\vv)$, $b(\vv) = \oint\vv\cdot\vm\dS$; the gauged
solution matches the Sherman--Morrison solve (the rank-one update formula
for the inverse) of the condensed problem to $2\times10^{-3}$ under a
degree-2 pressure, and the Lam\'e solution (the exact radially symmetric
response to uniform pressure) with a
fluid core to $3\times10^{-5}$ (2-D) and $5\times10^{-4}$ (3-D). Residual
decay $\sim0.1$ per step at $\eps = 10^{-2}$, $\mu_g = \kappa_f$;
observables agree to $10^{-4}$--$10^{-5}$ between $\eps = 10^{-2}$ and
$10^{-3}$ after refinement. Against the Dahlen path on a three-layer model
with an Adams--Williamson core: solid displacement and potential to
$2\times10^{-2}$ (the two discretisations' own error) under a zero-mean
load, the degree-0 difference present; Schur CG and block MINRES agree to
$10^{-6}$; serial and parallel to $10^{-6}$. Against \code{pyslfp} on
\code{prem\_4}: the gauged column of \S\ref{sec:cmb}.

%==========================================================================
\section{The fully referential self-gravitating problem}
\label{sec:referential}
%==========================================================================

The class \code{LinearQuasiStaticReferentialSelfGravitatingProblem} treats
a self-gravitating, arbitrarily pre-stressed elastic body about a general
equilibrium of a general reference body. Both the displacement and the
potential perturbation $\zeta^1$ are referential fields on the fixed
reference body; the physics of shape lives in the equilibrium mapping
$\varphi_e$.

\subsection{Three gravity formulations, and the choice}

With $\Phi_0$ the background potential:
(1) \emph{mixed, Eulerian potential} (WD eqs.~21--22; \S\ref{sec:mixed});
(2) \emph{mixed, referential potential} (Maitra \& Al-Attar
\cite{MaitraAlAttar2024}, eq.~21): $\zeta = \Phi\circ\varphi$ extended over
all space via a diffeomorphic extension of the motion, with Poisson operator
$\frac{1}{4\pi G}\langle a\grad\zeta,\grad\chi\rangle$,
\begin{equation}
  a(F) = J F^{-1}F^{-T};
\end{equation}
(3) \emph{eliminated, non-local} (AC18 eqs.~99/122): a dense double
integral, a theoretical device and never a good numerical formulation. The
class implements (2). It costs only assembly, leaves the DtN untouched (the
equilibrium mapping is the identity at and beyond the DtN sphere; the motion
of the exterior is gauge), and retires approximation (M5): the referential
form's gravity enters through $D\vu$ and composition alone, with no bare
displacement, and is indifferent to secular motion.

\subsection{The linearised referential system}

The static action \cite{MaitraAlAttar2024} is
\begin{equation}
  \mathcal{V}[\varphi,\zeta] = \int_M W(\vx,F)\dV + \int_M\rho\,\zeta\dV
  + \frac{1}{8\pi G}\int_{\R^3}\langle a(F)\grad\zeta,\grad\zeta\rangle\dV,
\end{equation}
with the motion extended diffeomorphically to $\R^3$. Perturb $\varphi =
\varphi_e + \vu$, $\zeta = \zeta^0 + \zeta^1$. With $H_u = F_e^{-1}D\vu$,
$a_e = a(F_e)$, and $a'(\vu)$, $a''(\vu,\vv)$ the first and (polarised)
second derivatives of $F\mapsto a(F)$ at $F_e$ in the directions $D\vu$,
$D\vv$,
\begin{align}
  a'(\vu) &= (\tr H_u)\,a_e - H_ua_e - a_eH_u^T, \\
  a''(\vu,\vv) &= [\tr H_u\tr H_v - \tr(H_uH_v)]\,a_e
     - \tr H_u(H_va_e + a_eH_v^T) - \tr H_v(H_ua_e + a_eH_u^T) \nonumber\\
  &\quad + (H_uH_v + H_vH_u)\,a_e + a_e(H_u^TH_v^T + H_v^TH_u^T)
     + H_ua_eH_v^T + H_va_eH_u^T ,
\end{align}
symmetric in $\vu\leftrightarrow\vv$. The density is a fixed referential
field and is never differentiated. With $\vg_0 = \grad\zeta^0$ and the
referential gravity flux $\vw_0 = a_e\vg_0$, the linearised system in
$(\vu,\zeta^1)$ with tests $(\vv,\chi)$ is
\begin{align}
  &a_{\mathrm{ref}}(\vu,\vv)
   + \frac{1}{8\pi G}\int\langle a''(\vu,\vv)\vg_0,\vg_0\rangle\dV
   + \frac{1}{4\pi G}\int\langle a'(\vv)\vg_0,\grad\zeta^1\rangle\dV
   = \ell_u(\vv), \label{eq:refu}\\
  &\frac{1}{4\pi G}\Bigl[\int_B\langle a_e\grad\zeta^1,\grad\chi\rangle\dV
     + \DtN(\zeta^1,\chi)\Bigr]
   + \frac{1}{4\pi G}\int\langle a'(\vu)\vg_0,\grad\chi\rangle\dV
   = \ell_\zeta(\chi), \label{eq:refz}
\end{align}
with $a_{\mathrm{ref}}$ the material + geometric stiffness \eqref{eq:matgeo}
and $\ell_u$, $\ell_\zeta$ the load forms (below).
The coupling expands to
$\langle a'(\vu)\vg_0,\grad\chi\rangle = (\tr H_u)(\vw_0\cdot\grad\chi)
- (H_u\vw_0)\cdot\grad\chi - (H_u^T\vg_0)\cdot(a_e\grad\chi)$, and the
gravity--gravity term reduces, for a vector basis function
$\phi_a\mathbf{e}_k$ (scalar basis function $\phi_a$ times the $k$-th
Cartesian unit vector)
(for which $H$ is rank one), to products of per-dof scalars built from
$F_e^{-T}\grad\phi_a$, $\grad\phi_a\cdot\vw_0$, $F_e^{-T}\vg_0$ and
$F_e^{-1}\vw_0$: $O(d)$ work per dof pair, no fourth-order coefficient
formed. The integrals of the gravity terms extend over the buffer as well
as the body (\S\ref{sec:vacuum}).

\paragraph{Structural remarks.}
\begin{itemize}
\item At $\varphi_e = \mathrm{id}$, hydrostatic: $H = D\vu$, $a_e =
  \mathbf{1}$, $\vw_0 = \vg_0 = \grad\Phi_0$, and the coupling is
  $\langle[(\divg\vu)\mathbf{1} - D\vu - D\vu^T]\grad\Phi_0,\grad\chi
  \rangle/4\pi G$. The system is related to the Eulerian one by the
  invertible change of variables
  \begin{equation}
    \zeta^1 = \phi + \vu\cdot\grad\Phi_0 :
    \label{eq:zetaphi}
  \end{equation}
  congruent, not identical, so cross-checks map solutions, not operators.
\item No gravity interface terms: $\rho$ is a fixed referential field, so
  no $[\grad\phi\cdot\vn + 4\pi G\rho\vu\cdot\vn]$-type terms and no
  $\grad\grad\Phi_0$ appear --- but see \S\ref{sec:vacuum}.
\item Saddle structure as before: symmetric, elastic-positive in $\vu$,
  Laplace-type in $\zeta^1$.
\item Loads: a surface mass load $\sigma$ (per referential area) loads the
  potential row alone, $\ell_\zeta(\chi) = -\int_{\partial M}\sigma\chi\dS$;
  the Eulerian $-\int\sigma\grad\Phi_0\cdot\vv$ is absorbed by
  \eqref{eq:zetaphi}.
\item The background $\zeta^0$ solves the referential Poisson equation with
  the fixed referential density (at construction, unless supplied).
\end{itemize}

\subsection{Assumptions and restrictions}

\begin{enumerate}
\item[(R1)] Linearised quasi-static, as throughout; rotation only through
  the centrifugal background.
\item[(R2)] General reference state $(\Ch,\Se,\varphi_e)$; $\Ch$ must be
  the \emph{bare} tensor (\S\ref{sec:dictionary}), and $\Se$ must be an
  equilibrium stress consistent with $\rho$ and the gravity (when generated
  by the background module it is; a raw coefficient is the modeller's
  responsibility).
\item[(R3)] The equilibrium mapping must extend over the whole ball as a
  diffeomorphism and be the identity at the DtN sphere (asserted, not
  enforced). The extension inside the buffer is gauge, but a rule is
  needed (\S\ref{sec:mappings}).
\item[(R4)] Fluid regions in this class are treated by the welded gauged
  formulation (\S\ref{sec:ref-gauged}); a genuinely slipping interface
  needs \S\ref{sec:slip}.
\end{enumerate}

\subsection{Background states}
\label{sec:background}

The general form takes $T^0$ (or $\Se$) as data, but it must satisfy
$\Div T^0 = \rho\grad\Phi_0$ with continuous tractions and be hydrostatic
in fluid regions. The background module provides:
\begin{itemize}
\item \code{RadialHydrostaticBackground}: hydrostatic $g$ and $p^0$ by
  cumulative quadrature for a radial model, the bare tensor from
  seismological $(\kappa,\mu)$ via \eqref{eq:bare}, $\Se = -p^0\mathbf{1}$,
  the identity mapping, and the assembled rheology;
\item \code{RelabelledBackground}: the same physical equilibrium from a
  relabelled reference body, owning the transformed chain
  $\tilde\rho = J_\xi\,\rho\circ\xi$,
  $\tilde C = J_\xi\,Q_\xi^{-T}(\Ch\circ\xi)\,Q_\xi^{-1}$ (Mandel form, with
  $Q_\xi$ the Mandel-form matrix of the strain transformation induced by
  $F_\xi$, as written in \code{background.hpp}),
  $\tilde S = J_\xi F_\xi^{-1}(\Se\circ\xi)F_\xi^{-T}$, $\varphi_e = \xi$,
  and the referential pressure $\tilde\pi = p^0\circ\xi$;
\item the Al-Attar \& Woodhouse \cite{AW10} generators for general
  aspherical $\Se$. Any two equilibrium stresses differ by an element of
  $\ker\Div_0$, the space of symmetric stress fields that are
  divergence-free in the body and traction-free on its surface. The minimum-norm
  equilibrium stress solves a linear elastic boundary-value problem
  (\code{MinimumNormEquilibriumStress}); the minimum-\emph{deviatoric} one
  solves a steady incompressible Stokes problem
  (\code{MinimumDeviatoricEquilibriumStress}), discretised with
  Taylor--Hood elements (pressure one order below velocity; equal order
  violates the inf--sup (Ladyzhenskaya--Babu\v{s}ka--Brezzi) stability
  condition of mixed velocity--pressure spaces, producing spurious pressure
  modes, and is refused). Both are matrix
  coefficients usable directly as $\Se$, and both admit relabelling (the
  generated $\Se$ is the second Piola--Kirchhoff pull-back
  $J F^{-1}(T\circ\varphi)F^{-T}$, generated on the fixed reference body).
  Solvability requires no net force or torque from body force and tractions
  (AW10 eq.~58), the compatibility enforced by the projected solvers.
\end{itemize}
Measured on a homogeneous ellipse of ellipticity $e$: loading responses with the full
generated $\Se$ versus its pressure part alone (the quasi-hydrostatic
approximation of standard practice, which is not an equilibrium stress off
sphericity) differ by $|\delta\vu|/|\vu|\approx0.13\,e$ at a pressure to
shear-modulus ratio $p/\mu = 0.31$,
scaling like the product $e\cdot(p/\mu)$; negligible at Earth's
non-hydrostatic figure, 1--7\,\% in displacement for $e = 0.1$--$0.3$.

\subsection{The mapping (relabelling) layer}
\label{sec:mappings}

\paragraph{Conventions.} $\xi: B\to\tilde B$ is a diffeomorphism of the
reference ball onto the physical domain, the identity on and outside the
DtN sphere (so $\code{SurfaceHarmonics}$, the DtN and boundary harmonic
analysis are unchanged; interface and surface topography live entirely in
$\xi$, the reference interfaces staying spheres so that the SubMesh coupling
is untouched). $F_{iA} = \partial\xi_i/\partial x_A$, $J = \det F > 0$.
Fields pull back componentwise ($u_i(\vx) = \tilde u_i(\xi(\vx))$; no
rotation of components). Three rules:
\begin{equation}
  \tilde\grad = F^{-T}\grad, \qquad
  \mathrm{d}\xi = J\dV, \qquad
  \tilde\vm\,\mathrm{d}\tilde S = \bnu\dS, \quad \bnu = JF^{-T}\vn
  \ \ \text{(Nanson)} .
\end{equation}
Integrators take the \emph{physical} material fields expressed in
reference coordinates ($\rho$ means $\tilde\rho\circ\xi$) with $J$, $F$
explicit in the form; nothing is absorbed into relabelled densities, so
with $\xi = \mathrm{id}$ every form reduces term-by-term.

\paragraph{The pulled-back forms.}
\begin{align*}
  &\text{Poisson:} && \textstyle\int J F^{-1}F^{-T}\grad\phi\cdot\grad\phi' \\
  &\text{elastic:} && \textstyle\int c'_{iAkB}\,\partial_Au_i\,\partial_Bu'_k,
       \quad c'_{iAkB} = J\,c_{ijkl}F^{-1}_{Aj}F^{-1}_{Bl} \\
  &\text{mass / advective:} && \textstyle\int J\rho\,\phi\phi', \quad
       \int J\rho\,(F^{-T}\grad\phi)\cdot\vu' \\
  &\text{divergence:} && \tilde\divg\tilde\vu = \tr(\grad\vu\,F^{-1}) \\
  &\text{surface:} && \textstyle\int_\Sigma f\,(\tilde\vm\cdot\vu)(\tilde\vm\cdot\vu')
       |\bnu|\dS,\quad \tilde\vm = \bnu/|\bnu| .
\end{align*}
$c'$ has only the major symmetry (45 components in 3-D) and is never
formed. \textbf{One assembly recipe} implements every domain term: at each
quadrature point replace $\mathrm{dshape}\to\mathrm{dshape}\cdot F^{-1}$
and the weight $w\to Jw$, then run the standard assembly (for the elastic
term: the usual Mandel $B$ matrix contracted with the referential
21-component tensor, exact because $\tilde c$ has its minor symmetries in
physical indices). Boundary terms use Nanson instead:
\code{BoundaryNormalScalarIntegrator} maps $(\vm\cdot\vv)\dS\to
(\bnu\cdot\vv)\dS$ (the $|\bnu|$ factors cancel), and
\code{BoundaryNormalNormalIntegrator} maps $(\vm\cdot\vu)(\vm\cdot\vu')\dS
\to(\bnu\cdot\vu)(\bnu\cdot\vu')/|\bnu|\dS$. Material symmetry survives:
an isotropic body pulls back to $c'_{iAkB} = J[\lambda G_{iA}G_{kB} +
\mu(\delta_{ik}(F^{-1}F^{-T})_{AB} + G_{kA}G_{iB})]$, $G = F^{-T}$
(the ``apparent anisotropy of a specific form'' of
\cite{AlAttarCrawford2016}, eq.~143).

\paragraph{Mapping objects.} A \code{Diffeomorphism} supplies $\xi$ and $F$
at a quadrature point: analytic (\code{CallableDiffeomorphism},
\code{RadialDiffeomorphism}, exact $F$) or discrete
(\code{GridFunctionDiffeomorphism}, $F = \mathbf{1} + \grad\mathbf{h}$). A
body-only mapping is extended to the ball either by
\code{TaperedDiffeomorphism} (blending an analytic ball-wide mapping to the
identity through a cubic smoothstep, $C^1$ at both seams, $F = \mathbf{1}$
at the DtN sphere) or by \code{NewHarmonicExtensionMapping} (nodal
interpolant on the body, extended into the buffer by a vector Laplace solve
with Dirichlet data on the body surface and zero on the DtN sphere).

\paragraph{Exact and interpolated modes; the change-of-variables
identity.} With exact $F$ the pulled-back solution converges to the mapped
exact solution (a convergence statement). With the nodal interpolant
$\xi_h$, assembling the standard integrators on the mapped mesh and the
mapped integrators on the reference mesh gives the \emph{same element
matrices to machine precision} (provided the same quadrature rule is used):
every mapped integrator is certified against its standard counterpart on an
arbitrarily deformed mesh. At solver level the identity holds to
$\sim10^{-6}$, bounded by the DtN centring on the shifted mesh centroid,
wherever every assembled term is covariant.

\paragraph{Shift maps require interpolated $F$.}
\label{sec:shiftmaps}
Mappings that move a fluid--solid interface (the \code{InterfaceShift}
maps of the perturbation benchmark, used at $\pm\eps$ for central-difference
derivatives with respect to interface position) are continuous but have a
\emph{kink}: their gradient $F$ jumps across the moving interface. The
one-sided slip-interface kernels (\S\ref{sec:slip}) evaluate $F$ exactly
there, and an exact, analytically branching map then returns the gradient
of whichever side the quadrature point falls on: on flat facets round-off
selects the side, while on curved facets the quadrature points sit $O(h^2)$
inside the analytic radius, so most of each facet systematically receives
the \emph{fluid}-side slope where the solid-side value belongs. This
map-evaluation defect was the entire outward-shift pathology of the AL
slip solver: with the exact map, the $+\eps$ shift ran $\sim9.6$k
iterations per degree with polluted Love numbers ($h'_1\approx-72$),
whereas with the per-SubMesh \emph{interpolated} map both shift signs show
identical, healthy behaviour ($\sim4$--$5$k iterations, sane Love numbers;
fluid-core model, $h = 0.3$). The interpolated map is side-consistent by
construction: each SubMesh's nodal interpolant samples only its own
elements, and the map's value is continuous. The benchmark therefore
forces interpolated $F$ for every shift run; the interpolation perturbs the
represented model by $O(h^p)$ identically at both signs, which cancels in
the central differences (\code{benchmarks/perturbation/README.md}, ``Shift
maps run with interpolated F (required)''). Two diagnostics of the broken
organisation used in that study,
\code{SetBrokenConstraintGravity} (supplying the continuous physical
gravity in place of the one-sided composition $\vb = F^{-T}\grad\zeta^0$
in the scalar-jump kernels) and \code{ScaleBrokenGravityInterface}
(scaling or dropping the $G_\Sigma$ block family, which changes the
physical problem), remain available for solver-structure comparisons only.

\paragraph{Equilibrium mapping versus relabelling pull-back.} Two distinct
uses of a mapping must not be confused. In \code{MaterialStiffnessIntegrator}
the mapping is the \emph{equilibrium mapping} $\varphi_e$ (the form
\eqref{eq:matgeo}); a relabelling composes into it and transforms the
coefficients by \eqref{eq:relabel} (owned by the background layer), leaving
the form unchanged. In the \code{Domain*}, \code{ElasticTensorIntegrator}
and boundary integrators the mapping is the \emph{relabelling pull-back}
with physical coefficients. The covariant fluid-gauge penalty must
therefore use \code{ElasticTensorIntegrator}$(C,\text{map})$, the certified
pull-back form, not \code{MaterialStiffnessIntegrator}, whose mapped form
expects a relabelled tensor.

\subsection{Discretisation: the vacuum extension}
\label{sec:vacuum}

The displacement lives on the body SubMesh, $\zeta^1$ on the ball. The
gravity terms of the second variation live wherever $\grad\zeta^0\neq0$,
the buffer included, through the \emph{extended} perturbation: as
$\varphi_e$ needs a taper, the linearised motion must be extended from
$\partial M$ to the DtN sphere, and $a'(\vu_{\mathrm{ext}})$,
$a''(\vu_{\mathrm{ext}},\cdot)$ contribute there. The extension is gauge
(observables are extension-independent; $\zeta^1$ itself is
extension-dependent in the vacuum), but it cannot be omitted: truncating
the gravity terms at $\partial M$ produces $O(1)$ errors. Options:
\begin{enumerate}
\item[(a)] \textbf{Ball-wide displacement} with a harmonic gauge penalty
  $\eps\mu_g\int\grad\vu:\grad\vv$ on the buffer and Tikhonov refinement
  (\code{SetVacuumExtension}). It works only in the \emph{biased} $O(\eps)$
  regime: the spectral condition of \S\ref{sec:tikhonov} fails, because the
  buffer has no physical stiffness (only the $L^2$-bounded zeroth-order
  gravity terms) while the harmonic penalty's Rayleigh quotients grow like
  $h^{-2}$; $\lambda\sim h^2\to0$ and refinement diverges with mesh or
  order. Kept as a findings log and for the non-linear problem, where the
  buffer motion becomes physical.
\item[(b)] \textbf{Prescribed linear extension} $\vu_{\mathrm{ext}} =
  E\,\vu|_{\partial M}$ with the buffer blocks folded through $E$ by sparse
  triple products (\code{SetPrescribedVacuumExtension}): exact (any $E$ is a
  gauge choice), no extra unknowns, no near-kernel, no $\eps$. \textbf{The
  route used.} \code{NewRadialVacuumExtension}: buffer dofs shared with the
  body copy their values exactly through the SubMesh dof pairing; interior
  buffer nodes at radius $r$ take $t(r)\,\vu(\vx_s)$ with
  \begin{equation}
    t(r) = \Bigl(\frac{r_{\mathrm{out}} - r}{r_{\mathrm{out}} - r_b}\Bigr)^2,
  \end{equation}
  so $t$ and $t'$ vanish at the DtN sphere (keeping $a = \mathbf{1}$
  there), $\vx_s$ the radial projection onto the body surface. In parallel
  the trace rows use the true-dof pairing and interior rows a query--reply
  exchange; the folds use hypre RAP/ParMult.
\end{enumerate}
Different tapers agree in observables: the built-in gauge-invariance test.
Note that the 2-D comparison with the Eulerian class must respect the
constant gauge ($\zeta^1$ and $\phi$ are both projected orthogonal to
constants, $\phi + \vu\cdot\grad\Phi_0$ is not).

\subsection{Solver and null space}

Block MINRES on the symmetric saddle system with the projected
block-diagonal preconditioner (AMG/Gauss--Seidel on the displacement block
and on the shifted \emph{mapped} Laplacian), serial and parallel in one
class, warm-started (\code{ResetSolution()} zeroes the iterate between
independent forcings, since the gauge refinement carries the previous
solution's fluid-gauge component forward).

\paragraph{Rigid modes have no potential partners.} $\zeta = \Phi\circ
\varphi$ is invariant under rigid motions of the body, so the null pairs are
$(\vt,0)$ for translations (every term vanishes identically; exact
discretely) and $(W\varphi_e,0)$ for rotations ($\sym(F_e^TD\vu) =
\sym(F_e^TWF_e) = 0$; geometric and gravity terms cancel by rotational
invariance of the equilibrium energy; near-null discretely). Contrast the
coupled pairs of \S\ref{sec:mixed-null}. Both are projected from the
displacement in the true-dof inner product; the potential needs no
projection (except the 2-D constant, as in \S\ref{sec:mixed-2d}).
\code{NullPairResidual} evaluates any candidate pair $(\vu,0)$ under the
full block operator.

\subsection{Gauged fluids in the referential class}
\label{sec:ref-gauged}

The welded gauged treatment transfers with no new machinery: the fluid's
inputs come from the background module (bare tensor with $\mu_b = p^0$,
$\Se = -p^0\mathbf{1}$), and the base-class penalty and refinement act on
the fluid attributes unchanged, with \code{GaugeRefine()} overridden as in
\S\ref{sec:tikhonov} (zero potential load in refinement solves, potential
accumulated). The deviatoric penalty is assembled covariantly through the
rheology's equilibrium mapping by default: since the gauge subspace has
stiffness $O(\eps)$, an $O(\eps A)$ difference in a non-covariant penalty
produces $O(A)$ differences in the gauge representative.

\paragraph{Verification.} (i) The joint relabelling-invariance identity: an
azimuthal relabelling $\vw = \mathrm{curl}\,\psi$ supported in the core is
annihilated by the displacement-row operator only \emph{jointly} --- the
material $\mu_b = p^0$ term, the geometric $-p^0$ term and the gravity
second variation cancelling through the equilibrium condition. With the
bare dictionary the normalised energy along the orbit converges to zero
(0.087, 0.0088, 0.0016 at orders 1--3); with seismological moduli used
directly it converges to the finite defect $\int2p^0|\sym D\vw|^2\approx
0.137$: an $85\times$ separation, the sharpest single test of
\eqref{eq:bare}. (ii) Against the Eulerian gauged class on a two-layer
fluid-core problem: agreement in mantle displacement and in
$\phi + \vu\cdot\grad\Phi_0$ at the 2\,\% discretisation level on the coarse
test mesh. (iii) The hydrostatic cross-check against the mixed class (the
``tier-(i)'' check: same physical problem, solutions mapped through
\eqref{eq:zetaphi}) with the prescribed
vacuum extension: displacement and potential map onto the Eulerian
solution at discretisation level, improving with order. (iv) Relabelled
welded-gauged solver identity: strict, $\vu$ to $1.6\times10^{-7}$ and
$\zeta$ to $4\times10^{-8}$ at relabelling amplitude $0.02$ (the size of
$\xi - \mathrm{id}$ relative to the body radius), element size $h = 0.3$.

%==========================================================================
\section{The slipping-interface problem}
\label{sec:slip}
%==========================================================================

The class \code{LinearQuasiStaticReferentialSelfGravitatingSlipProblem}
carries a broken displacement pair $(\vv_{\mathrm{s}},\vv_{\mathrm{f}})$
on solid and fluid SubMeshes of the ball with a genuinely slipping
fluid--solid interface, in the referential formulation
(\code{doc/slip\_interface.tex}). Following that note, $\vv$ denotes the
broken displacement, $M_{\mathrm{s}}$, $M_{\mathrm{f}}$ the reference solid
and fluid regions, $\vN$ the \emph{referential} unit normal on $\Sigma$
pointing \emph{out of the fluid}, and $\dS$ the referential measure.

\subsection{When it is needed; assumptions}

For a barotropic fluid the welded gauge is admissible and nothing here
arises. The broken formulation is forced for two-parameter (stratified)
fluids, where tangential continuity is not an admissible gauge, and for the
eventual non-linear slip map; when the broken space is used on a
pressurised background its interface terms are obligatory. Restrictions of
the implementation: (S1) one fluid region inside a solid shell
(nested-shell, two-sided extensions are not implemented); (S2) the
single-valued organisation is assembled at $\varphi_e = \mathrm{id}$ only
and refuses a non-identity mapping (\code{Diffeomorphism::IsIdentity}); the
broken-$\zeta$ organisation is assembled mapped throughout and is the route
for mapped and aspherical backgrounds; (S3) the fluid displacement is
gauge (linearised relabellings), fixed by the penalty and refinement of
\S\ref{sec:tikhonov} in both enforcement paths; (S4) the KKT path is
implemented for the single-valued organisation only.

\subsection{Kinematics and the constraint}

With $\varphi_e$ on the ball, $F_e$, $J_e$, the unnormalised \textbf{Nanson
normal}
\begin{equation}
  \bnu := \cof(F_e)\vN = J_eF_e^{-T}\vN, \qquad
  \vn\,\mathrm{d}S_{\mathrm{phys}} = \bnu\dS, \qquad
  |\bnu|\dS = \mathrm{d}S_{\mathrm{phys}} .
\end{equation}
The two conventions ``unit physical normal $\times$ physical measure'' and
``Nanson normal $\times$ referential measure'' give the same pairing; the
error mode is mixing them ($\vn$ against $\dS$ silently loses $|\bnu|$).
Every surface integral carries $(\bnu,\dS)$ explicitly.

The fluid stays attached but slips: $\varphi_{\mathrm{f}}|_\Sigma =
\varphi_{\mathrm{s}}|_\Sigma\circ\sigma$, with slip map $\sigma_\eps =
\mathrm{id} + \eps\vs + \tfrac{\eps^2}{2}\bm{\sigma}_2$, $\vs\in T\Sigma$.
The second-order coefficient has a forced normal part,
$\bm{\sigma}_2 = \vs_2 - \mathrm{II}(\vs,\vs)\vN$ ($\mathrm{II}$ the
referential second fundamental form). To first order $\vv_{\mathrm{f}}
|_\Sigma = \vv_{\mathrm{s}} + F_e\vs$, so with $\jump{\vv} :=
\vv_{\mathrm{s}} - \vv_{\mathrm{f}} = -F_e\vs$,
\begin{equation}
  \boxed{\ \bnu\cdot\jump{\vv} = 0\ }, \qquad \vs = -F_e^{-1}\jump{\vv},
  \label{eq:constraint}
\end{equation}
the normal-jump constraint with the \emph{mapped} normal. To second order
the curvature enters and re-hides: by the Weingarten identity
$\vN\cdot D\vs[\vs] = -\mathrm{II}(\vs,\vs)$ (Valette's content
\cite{Valette1986,Valette1991}), the $\bnu$-component of the fluid's
second-order boundary datum is
\begin{equation}
  \bnu\cdot\mathbf{b}_2 = \bnu\cdot\grad_\Sigma(\vv_{\mathrm{s}}
     + \vv_{\mathrm{f}})[\vs],
\end{equation}
tangential derivatives of the two traces only: neither curvature nor
$D^2\varphi_e$ is ever assembled.

\subsection{First variation: equilibrium and the multiplier}

With the fluid stored energy $V(\vx,J)$ and the \emph{referential pressure}
$\pi := -D_JV(\vx,J_e)$ ($= p^0\circ\varphi_e$), the fluid boundary
traction per referential area is $P_{\mathrm{f}}\vN = -\pi\bnu$.
Stationarity under constrained variations gives interior equilibrium and
the \textbf{equilibrium consistency condition}
\begin{equation}
  P_{\mathrm{s}}\vN = P_{\mathrm{f}}\vN = -\pi\bnu \quad\text{on }\Sigma,
\end{equation}
a checkable property of a generated background, and identifies the
multiplier of $\oint\lambda\,\bnu\cdot\jump{\vv}\dS$ as $\lambda = \pi$. In
the \emph{linearised} solve, however, the converged multiplier is the
first-order interface traction $\varpi^{(1)}$ (the multiplier of a
\emph{mortar}, i.e.\ Lagrange-multiplier, coupling of the two traces), not
$\pi$.

\subsection{Second variation: the elastic interface form}

The second variation along the constrained family is $\delta^2E +
\delta E[(\mathbf{0},\mathbf{u}_2)]$; the second term collapses by interior
equilibrium to $-\oint\pi\,\bnu\cdot\mathbf{b}_2\dS$, giving:
\begin{proposition}[Interface form]
The entire elastic interface contribution is
\begin{equation}
  B_\Sigma(\vv,\vv') = -\tfrac12\oint_\Sigma\pi\Bigl[
     \bnu\cdot\grad_\Sigma(\vv_{\mathrm{s}}+\vv_{\mathrm{f}})[\vs']
   + \bnu\cdot\grad_\Sigma(\vv'_{\mathrm{s}}+\vv'_{\mathrm{f}})[\vs]\Bigr]\dS,
  \qquad \vs = -F_e^{-1}\jump{\vv},
  \label{eq:BSigma}
\end{equation}
the mean-trace derivative against the jump; there are no other elastic
interface terms.
\end{proposition}
The fluid volume Hessian, from $V(\vx,J)$ with $H = F_e^{-1}D\vv_{\mathrm{f}}$,
is $\int[V_{JJ}J_e^2(\tr H)^2 - \pi J_e((\tr H)^2 - \tr H^2)]$, which equals
the general material + geometric split with the \emph{bare} fluid moduli:
$\mu_b = \pi$ follows from the fluid's constitutive class alone,
independently confirming the dictionary \eqref{eq:bare}.

\subsection{Gravity, single-valued organisation}

Let $\tilde\varphi$ be the solid motion extended smoothly over
$M_{\mathrm{f}}$ (and the buffer), $\zeta := \Phi\circ\tilde\varphi$
single-valued on the ball (ordinary $H^1$, DtN untouched), and $\vw :=
\vv_{\mathrm{f}} - \tilde\vv$ the fluid-versus-extension mismatch
($\bnu\cdot\vw|_\Sigma = 0$ by \eqref{eq:constraint}). The fluid source
composes through $\tilde\varphi^{-1}\circ\varphi_{\mathrm{f}} = \mathrm{id}
+ \eps F_e^{-1}\vw + \dots$.
\begin{lemma}[No gravity interface terms]
At first order the boundary term of the composed source is
$\oint\rho_{\mathrm{f}}(\vw\cdot\bnu/J_e)\zeta\dS = 0$; at second order the
by-parts boundary term is again proportional to $\bnu\cdot\vw = 0$. The
slip's gravitational physics enters only through volume terms on
$M_{\mathrm{f}}$.
\end{lemma}
With $\vb := F_e^{-T}\grad\zeta^0$, the additions to the operators of
\S\ref{sec:referential} (acting with the extended field ball-wide) are, in
Hessian normalisation,
\begin{align}
  H_{\zeta\vw} &= 2\int_{M_{\mathrm{f}}}\rho\,(F_e^{-1}\vw)\cdot\grad\zeta^1\dV
     \qquad\text{(mismatch coupling)}, \\
  H_{\vw\vw} &= -\int_{M_{\mathrm{f}}}\Div(\rho F_e^{-1}\vw)\,
     \bigl((F_e^{-1}\vw)\cdot\grad\zeta^0\bigr)
   - \int_{M_{\mathrm{f}}}\rho\,\vb\cdot D\vw[F_e^{-1}\vw]
   - 2\int_{M_{\mathrm{f}}}\rho\,\vb\cdot D\tilde\vv[F_e^{-1}\vw],
\end{align}
all first-derivative volume forms (a $D^2\varphi_e$ term cancels
identically, the volume analogue of the curvature collapse). Eliminating
$\zeta^1$ at fixed motion contributes a negative-definite term
($-\frac{1}{4\pi G}\int|\grad\zeta^{1*}|^2$): self-gravity destabilises.

\subsection{\texorpdfstring{Gravity, broken-$\zeta$ organisation}{Gravity, broken-zeta organisation}}

Alternatively compose region-wise with each region's own motion, $\zeta_r
= \Phi\circ\varphi_r$ ($r$ = fluid; outer = solid with its vacuum
continuation, whose buffer extension survives as gauge). The field energy
splits over the deformed regions (a tiling), giving
\begin{equation}
  E_g = \operatorname{stat}\Bigl[\sum_r\frac{1}{8\pi G}\int_{M_r}
    \langle a(F_r)\grad\zeta_r,\grad\zeta_r\rangle
    + \sum_r\int_{M_r}\rho\,\zeta_r\Bigr]
\end{equation}
subject to $\zeta_{\mathrm{f}} = \zeta_{\mathrm{s}}\circ\sigma$ on
$\Sigma$. Both sources are exact: the fluid extension, $\vw$ and the
mismatch terms disappear identically. Linearising,
\begin{equation}
  \zeta^1_r = \varphi^1 + \vb\cdot\vv_r \quad\Longrightarrow\quad
  \boxed{\ \jump{\zeta^1} = \vb\cdot\jump{\vv} = -\vs\cdot\grad_\Sigma\zeta^0\ },
  \label{eq:zjump}
\end{equation}
imposed in the first form: a homogeneous linear constraint needing neither
tangential projection nor the normal-jump constraint. With the background
flux density $q := \frac{1}{4\pi G}\vb\cdot\bnu$:
\begin{proposition}[Gravity interface form]
\begin{equation}
  G_\Sigma = \oint_\Sigma\Bigl[\frac{|\vb|^2}{8\pi G}\,\bnu\cdot
     \grad_\Sigma(\vv_{\mathrm{s}}+\vv_{\mathrm{f}})[\vs]
   + q\Bigl(\grad_\Sigma(\zeta^1_{\mathrm{s}}+\zeta^1_{\mathrm{f}})[\vs]
     - \vb\cdot\grad_\Sigma(\vv_{\mathrm{s}}+\vv_{\mathrm{f}})[\vs]\Bigr)
   \Bigr]\dS,
\end{equation}
used in symmetrised polarisation; coefficients $\bnu$, $\vb$, $q$ and
tangential trace derivatives only.
\end{proposition}
The collected broken operator: unknowns $(\vv_{\mathrm{s}},\vv_{\mathrm{f}},
\zeta^1_{\mathrm{o}},\zeta^1_{\mathrm{f}})$; per region the operators of
\S\ref{sec:referential} with each region's own field (mapped Poisson block
per region, DtN on the outer one); per interface $B_\Sigma$, $G_\Sigma$ and
the two trace constraints $\bnu\cdot\jump{\vv} = 0$ and \eqref{eq:zjump}.
Limits: welded ($\vs\equiv0$) gives the continuous-$\zeta$ welded problem;
gravity off ($\vb = 0$) gives the sliding tests. Everything is
per-interface, so topology scales trivially. The two organisations are
exactly equivalent (a bijective change of variables) and share \emph{no}
gravity-interface machinery, which makes their head-to-head a two-sided
check.

\subsection{Discretisation}

\paragraph{Two spaces and the nodal pairing.} $\vv_{\mathrm{s}}$ and
$\vv_{\mathrm{f}}$ live in independent $H^1$ spaces on two SubMeshes of one
parent; the only broken surface is $\Sigma$ (no interior penalty, no
skeleton terms). Each space has a signed dof injection into the parent
space ($\mathcal{I}$ of \S\ref{sec:submesh}, written $\Pi_{\mathrm{s}}$,
$\Pi_{\mathrm{f}}$ here), and
\begin{equation}
  J = \Pi_{\mathrm{s}}^T\Pi_{\mathrm{f}}
  \qquad(\text{solid}\times\text{fluid dofs, entries }\pm1)
\end{equation}
identifies the interface dofs, orientation included, with no geometric
matching (in parallel a hypre product of the true-dof injections, which
handles interface dofs whose two sides are owned by different processes).
Because the two trace spaces are the same discrete space in different
coordinates, every interface form is assembled \textbf{once}, on the solid
side's boundary elements. Strong nodal imposition of the normal constraint
is not used: the discrete normal is ill-defined at the vertices of a
faceted interface where the dofs sit; the weak form evaluates normals only
at face-interior quadrature points.

\paragraph{Interface blocks.} Let $U = (U_{\mathrm{s}}, U_{\mathrm{f}})$
be the coefficient vectors of the broken pair, $B_n$ the matrix of the
(mapped) normal--normal boundary form $\int_\Sigma(\bnu\cdot\vu)(\bnu\cdot
\vu')/|\bnu|\dS$ on the solid side, $\tilde\vm = \bnu/|\bnu|$, and
$\mathsf{D} = [I, -J]$ the jump map ($\mathsf{D}U = U_{\mathrm{s}} -
JU_{\mathrm{f}}$, the discrete jump on solid-side dofs). With a penalty
parameter $\theta > 0$ the normal-jump penalty is
\begin{equation}
  \theta P,\qquad P = \mathsf{D}^TB_n\mathsf{D}
  = \begin{pmatrix} B_n & -B_nJ\\ -J^TB_n & J^TB_nJ\end{pmatrix},
  \qquad \theta\,U^TPU = \theta\int_\Sigma\bigl((\vv_{\mathrm{s}}-\vv_{\mathrm{f}})
     \cdot\tilde{\vm}\bigr)^2|\bnu|\dS ,
\end{equation}
positive semi-definite with kernel all fields of continuous normal trace
(every tangential jump). $B_\Sigma$ is formed from the one-sided kernel
$G_\pi$ of \code{SlipInterfacePressureIntegrator} (assembled on the solid
side with coefficient $\pi$) with the sum map $\mathsf{S} = [I, J]$ and the
jump map $\mathsf{D}$ as $B_\Sigma = -\frac12(\mathsf{S}^TG_\pi\mathsf{D} +
\mathsf{D}^TG_\pi^T\mathsf{S})$ (the minus because $G_\pi$ uses the solid's
outward normal, $-\vN$). $G_\Sigma$ is $\frac12(\mathsf{S}^TG_A\mathsf{D} +
\mathsf{D}^TG_A^T\mathsf{S}) + [\frac12\mathsf{D}^TM_q\mathsf{S}_\zeta +
\text{transpose}]$, with $G_A$ and $M_q$ the one-sided vector and
scalar--vector kernels (\code{SlipInterfaceGravityIntegrator},
\code{SlipInterfaceGravityScalarIntegrator}), $J_\zeta$ the scalar dof
pairing and $\mathsf{S}_\zeta = [I, J_\zeta]$ (signs plus: both $\bnu$-linear coefficients flip
with the normal, cancelling the minus of $\vs = -F_e^{-1}\jump{\vv}$).
These sign conventions are pinned by discrete-versus-quadrature
cross-checks. All these one-sided kernels evaluate $F_e$ on the interface
itself, which is why an equilibrium mapping with a gradient kink at
$\Sigma$ (an interface-shift map) must be supplied in the per-SubMesh
interpolated form (\S\ref{sec:shiftmaps}).

\paragraph{Single-valued gravity pieces.} The fluid extension $\tilde\vv =
E_f\vv_{\mathrm{s}}$ is a prescribed sparse operator
(\code{NewRadialFluidExtension}): interface trace rows copy the solid
values exactly (to round-off, the property the vanishing lemmas need);
interior rows take $t(r)\vv_{\mathrm{s}}(r_c\hat\vx)$, with $r_c$ the
interface radius, $\hat\vx = \vx/|\vx|$ and $t = (r/r_c)^{p_t}$ vanishing at
the centre (taper power $p_t = 2$ by default); the interior rule is gauge. At
$\varphi_e = \mathrm{id}$ the volume pieces assemble from library
integrators: $\int\rho\,\vw\cdot\grad\grad\zeta^0\,\vw'$ as a vector mass
with matrix coefficient; the mismatch coupling as a mixed
vector--scalar-gradient form; and the $\tilde\vv$-term as the
\emph{compensated} combination $G_{\mathrm{vg}} - M_{\mathrm{H}}$, with
$G_{\mathrm{vg}}$ the library's vector--gradient--vector form and
$M_{\mathrm{H}}$ the matrix-coefficient mass term just described: the
former differentiates the nodally interpolated product
$\grad\zeta^0\cdot\tilde\vv$ and so includes exactly the mass term
already assembled. The solid row receives the fluid-region gravity folded
through $E_f$ and the mismatch folds.

\paragraph{Broken-$\zeta$ pieces.} Per-region potentials $(\zeta_o,
\zeta_f)$ on the outer (solid + buffer) and fluid SubMeshes; the DtN stays
on the ball space and is folded through the outer injection; the
scalar-jump constraint uses the solid-side kernels $M_z$ (scalar boundary
mass), $K_{vz} = \oint(\vb\cdot\vu)\zeta$ and $P_b$ (the $\vb\otimes\vb$
vector boundary mass), assembled through the same one-sided pairings. The
vacuum extension of \S\ref{sec:vacuum} is still required.

\subsection{Constraint enforcement I: penalty and augmented Lagrangian}
\label{sec:al}

Collect the discrete single-valued problem as follows: $U = (U_{\mathrm{s}},
U_{\mathrm{f}}, Z)$ with $Z$ the potential coefficients; $\mathcal{A}$ the
symmetric physical block operator (elastic, geometric, gravity, mismatch
and $B_\Sigma$ blocks); $f$ the load; $\eps Q$ the fluid-gauge penalty
(\S\ref{sec:tikhonov}) on the $U_{\mathrm{f}}$ block; and $P$ the
constraint penalty above, acting on the displacement blocks. The physical
problem is: find $U$ and a multiplier (a dual vector) $w^*\in
\mathrm{range}(P)$ with
\begin{equation}
  \mathcal{A}U + w^* = f, \qquad PU = 0 .
  \label{eq:constrained}
\end{equation}
$w^*$ is the discrete interface normal-traction perturbation paired with
$\bnu$. A pure penalty, $(\mathcal{A} + \theta P)U = f$, enforces $PU = 0$
only as $\theta\to\infty$, where $\mathcal{A} + \theta P$ becomes
arbitrarily ill-conditioned.

\paragraph{The multiplier iteration.} The augmented-Lagrangian (Uzawa-type)
iteration keeps $\theta$ moderate and updates a running multiplier $w$:
\begin{equation}
  (\mathcal{A} + \theta P)\,U^{(k+1)} = f - w^{(k)}, \qquad
  w^{(k+1)} = w^{(k)} + \theta P\,U^{(k+1)}, \qquad w^{(0)} = 0 .
  \label{eq:alplain}
\end{equation}
At a fixed point the update stalls, so $PU = 0$, and then
$\mathcal{A}U + w = f$: the fixed point is \eqref{eq:constrained}, exact at
finite $\theta$ and independent of it.

\paragraph{Contraction rate.} Let $(U^*, w^*)$ solve
\eqref{eq:constrained}; since $PU^* = 0$, $(\mathcal{A} + \theta P)U^* = f -
w^*$. Subtracting, the errors $e_U^{(k)} = U^{(k)} - U^*$,
$e_w^{(k)} = w^{(k)} - w^*$ obey
\[
  (\mathcal{A} + \theta P)\,e_U^{(k+1)} = -e_w^{(k)}, \qquad
  e_w^{(k+1)} = e_w^{(k)} + \theta P e_U^{(k+1)}
  = \bigl(I - \theta P(\mathcal{A} + \theta P)^{-1}\bigr)e_w^{(k)} .
\]
For the analysis factor $P = R^TR$ (e.g.\ $R = B_n^{1/2}\mathsf{D}$; never
formed), assume $\mathcal{A}$ invertible (on the complement of the
projected null pairs), and write $e_w = R^T\mu$ (all multipliers lie in
$\mathrm{range}(P) = \mathrm{range}(R^T)$). The Sherman--Morrison--Woodbury
identity gives, with the constraint-space Schur complement
$\hat S := R\mathcal{A}^{-1}R^T$,
\begin{equation}
  R(\mathcal{A} + \theta R^TR)^{-1}R^T = \hat S(I + \theta\hat S)^{-1},
  \label{eq:woodbury}
\end{equation}
hence
\begin{equation}
  e_w^{(k+1)} = R^T\bigl[I - \theta\hat S(I + \theta\hat S)^{-1}\bigr]\mu^{(k)}
             = R^T(I + \theta\hat S)^{-1}\mu^{(k)} .
\end{equation}
On an eigenvector of $\hat S$ with eigenvalue $\sigma_j$ the multiplier
error is multiplied per sweep by
\begin{equation}
  \frac{1}{1 + \theta\sigma_j} = \frac{\sigma_j^{-1}}{\sigma_j^{-1} + \theta}
  \;\le\; \frac{\norm{A_S}}{\norm{A_S} + \theta},
  \qquad A_S := \hat S^{-1},
\end{equation}
the estimate of the notes: $A_S$ is the stiffness that the constraint
``sees'', and the slowest mode is the constraint direction along which
$\mathcal{A}$ is stiffest. The constraint violation
$RU^{(k+1)} = Re_U^{(k+1)} = -\hat S(I + \theta\hat S)^{-1}\mu^{(k)}$
contracts at the same rate. Larger $\theta$ contracts faster per sweep but
worsens the conditioning of $\mathcal{A} + \theta P$, i.e.\ the cost of each
sweep; the measured compromise is $\theta = 100$, below which a
conditioning cliff appears, above which a flat plateau.

\paragraph{Interleaving with the gauge refinement.} In the slip class the
fluid-gauge refinement of \S\ref{sec:tikhonov} and the multiplier update
run in one loop,
\begin{equation}
  (\mathcal{A} + \eps Q + \theta P)\,U^{(k+1)} = f - w^{(k)} + \eps Q\,U^{(k)},
  \qquad w^{(k+1)} = w^{(k)} + \theta P\,U^{(k+1)} .
  \label{eq:al}
\end{equation}
The added $\eps QU^{(k)}$ is the Tikhonov correction: at a fixed point it
cancels the penalty, $PU = 0$, and $U$ solves \eqref{eq:constrained} with
neither $\theta$ nor $\eps$ in it. Each part contracts at its own rate in
isolation; the combined loop is observed to contract the normal jump by
$\sim10^2$ over the sweeps, to a floor set by the interleaved gauge source
($\sim6\times10^{-5}$ of the normal trace on the coarse test disc,
$\theta$-independent), while the tangential jump remains $O(1)$.

\paragraph{Algorithm (AL path, \code{SetConstraint}$(\theta, K)$).}
\begin{enumerate}
\item Assemble $\mathcal{A}$, $\eps Q$, $P$; build the solver blocks
  $\mathcal{A} + \eps Q + \theta P$ and the block-diagonal preconditioner
  (the problem's own preconditioners rebuilt on the penalised blocks: AMG
  with systems options on the fluid block in parallel, GS serially; the
  shifted mapped Laplacian on $Z$); build the projector of the null pairs
  (\S\ref{sec:slip-null}).
\item Set $w\leftarrow0$ and take $U$ from the warm start.
\item For $k = 0,\dots,K-1$ ($K$ = \code{al\_iterations}, default 8):
  \begin{enumerate}
  \item form the right-hand side $f - w + \eps QU$;
  \item solve with projected block MINRES (\S\ref{sec:projected}) to the
    sweep tolerance $\tau_k$ (below), giving the new $U$;
  \item update $w\leftarrow w + \theta PU$;
  \item record the normal-jump energy $(j^TB_nj)^{1/2}$, $j = U_{\mathrm{s}}
    - JU_{\mathrm{f}}$ (\code{NormalJumpHistory()}).
  \end{enumerate}
\item Return $U$ and $w$ (\code{ConstraintMultiplier()}).
\end{enumerate}
In the broken-$\zeta$ organisation the scalar-jump constraint
\eqref{eq:zjump} receives its own penalty $\theta_\zeta P_\zeta$, with
$U^TP_\zeta U = \oint_\Sigma(\jump{\zeta^1} - \vb\cdot\jump{\vv})^2\dS$
assembled from the kernels $M_z$, $K_{vz}$, $P_b$, and its own multiplier
updated in step 3(c) alongside $w$; the four-block system $(U_{\mathrm{s}},
U_{\mathrm{f}}, Z_o, Z_f)$ is solved by MINRES in the same way, and the
scalar-jump energies are recorded (\code{ZetaJumpHistory()}).

\paragraph{Inexact sweeps.} Only the last sweep determines the endpoint
to full accuracy: an error $\eta$ left by an inexact solve in sweep $k$
enters $w$ as $\theta P\eta$ and is then contracted by the remaining sweeps
like any multiplier error. \code{SetSweepTolerance}$(\tau_{\mathrm{loose}})$
exploits this: the early sweeps run at a relaxed relative tolerance,
tightened geometrically from $\tau_{\mathrm{loose}}$ towards the solver's
tolerance $\tau$ (for instance $\tau_k = \tau_{\mathrm{loose}}
(\tau/\tau_{\mathrm{loose}})^{k/(K-1)}$), and the final sweep always runs
at full tolerance. The same schedule applies to the KKT path's gauge
refinements (\S\ref{sec:kkt}). Measured endpoint shifts are $\sim10^{-3}$
(AL) and $\sim2\times10^{-2}$ (KKT) of the spread between the exact
endpoints, an order below mesh error, at roughly half the cost. $\tau_{\mathrm{loose}} = 0$ (the default) keeps
every sweep at full tolerance --- the reproducible-endpoint mode that
finite-difference studies and strict comparisons require; the benchmark
enforces it on interface-shift runs (\S\ref{sec:shiftmaps}). Two schedules
failed and are not used: an inner tolerance coupled to the measured jump
deadlocks (the measured jump cannot fall below the floor the loose
tolerance itself sets, while multiplier errors are amplified by $\theta$),
and a plateau-based early exit false-triggers on the same measurement
noise and truncates the multiplier updates.

\subsection{Constraint enforcement II: the KKT system}
\label{sec:kkt}

\paragraph{The system.} The KKT path (\code{EnableKKT}) makes the
multiplier an unknown: $\lambda$, a coefficient vector on the interface
dofs of a scalar space on the solid SubMesh. The constraint row is
$\mathcal{C}U = N^T(U_{\mathrm{s}} - JU_{\mathrm{f}})$ restricted to those
dofs, with $N$ the one-sided flux kernel $\oint_\Sigma\lambda\,
(\bnu\cdot\vv)\dS$ (\code{BoundaryNormalScalarIntegrator} with the
equilibrium mapping; Nanson's relation \eqref{eq:nanson0} makes such
flux-type forms exact under mappings, so the row is covariant by
construction). With $\hat M = \mathrm{diag}(\hat m_i)$ the lumped interface
mass of the multiplier space, one MINRES solves
\begin{equation}
  \begin{pmatrix} S_\theta & \mathcal{C}^T \\ \mathcal{C} & 0\end{pmatrix}
  \begin{pmatrix} U\\ \lambda\end{pmatrix} =
  \begin{pmatrix} f\\ 0\end{pmatrix},
  \qquad
  S_\theta := \mathcal{A}_\eps + \theta\,\mathcal{C}^T\hat M^{-1}\mathcal{C},
  \quad \mathcal{A}_\eps := \mathcal{A} + \eps Q .
  \label{eq:kkt}
\end{equation}
The outer loop reduces to the fluid-gauge refinements of
\S\ref{sec:tikhonov}: the gauge indeterminacy of $U_{\mathrm{f}}$ is an
underdetermination, not an interface condition, and stays penalty +
refinement in both paths.

\paragraph{The augmentation does not change the solution.} Any solution
of \eqref{eq:kkt} has $\mathcal{C}U = 0$, so the augmentation term
$\theta\mathcal{C}^T\hat M^{-1}\mathcal{C}U$ vanishes and $(U,\lambda)$
solves the unaugmented system, for every $\theta$. This is why the KKT
endpoint is invariant to $\theta\in[1,300]$ at the $5\times10^{-10}$ level
(whereas the AL endpoint moves with $\theta$ through its constraint floor)
--- a sharp regression check on the constraint blocks. $\theta$ affects
only the conditioning.

\paragraph{Why the Schur complement clusters at $\hat M/\theta$
(Golub--Greif).} Eliminating $U$ from \eqref{eq:kkt} gives the multiplier
Schur complement $\mathcal{S} = \mathcal{C}S_\theta^{-1}\mathcal{C}^T$.
Write $R = \hat M^{-1/2}\mathcal{C}$, so that the augmentation is
$\theta R^TR$, and $\hat S = R\mathcal{A}_\eps^{-1}R^T$. The Woodbury
identity \eqref{eq:woodbury} gives
\begin{equation}
  \hat M^{-1/2}\,\mathcal{S}\,\hat M^{-1/2} = R S_\theta^{-1}R^T
  = \hat S(I + \theta\hat S)^{-1},
\end{equation}
whose eigenvalues are $\sigma_j/(1 + \theta\sigma_j)\in(0, 1/\theta)$,
$\sigma_j$ those of $\hat S$. Equivalently, the preconditioned Schur
complement $(\theta\hat M^{-1})\,\mathcal{S}$ is similar to
$\theta\hat S(I + \theta\hat S)^{-1}$, with eigenvalues
\begin{equation}
  \frac{\theta\sigma_j}{1 + \theta\sigma_j}\in(0,1),
  \qquad \to 1 \ \text{ as } \theta\sigma_j\to\infty
\end{equation}
(and exactly 1 on constraint directions where $\mathcal{A}_\eps$ is
singular, the Golub--Greif setting \cite{GolubGreif2003}). Constraint
directions on which $\mathcal{A}_\eps$ is soft ($\sigma_j$ large) are
therefore clustered at 1; only directions with $\theta\sigma_j\lesssim1$
spread the spectrum, and raising $\theta$ moves them towards 1 ---
consistent with the measured iteration counts ($\theta = 100$: 16.9k;
10: 27.7k; 1: 49.5k) even though the multiplier, not the penalty, enforces
the constraint. The argument needs the augmentation to be built from
\emph{the same} $\mathcal{C}$ and $\hat M$ as the constraint row and its
preconditioner. The AL penalty $\theta P = \theta\mathsf{D}^TB_n\mathsf{D}$
(the full $L^2$ trace form, not its projection onto the multiplier space)
would leave a kernel mismatch of the size of the gap between the
equal-order trace space and the multiplier space, and the identity above
would not hold; replacing it by the consistent
$\theta\mathcal{C}^T\hat M^{-1}\mathcal{C}$ was worth $1.35\times$ (22.9k
$\to$ 16.9k iterations). It is assembled like the $B_n$ folds, as
$N\,\mathrm{diag}(1/\hat m_i)\,N^T$ expanded over $(U_{\mathrm{s}},
U_{\mathrm{f}})$ with $J$ (serial sparse products or hypre products).

\paragraph{The preconditioner.} Block-diagonal, symmetric positive
definite, as MINRES requires:
\begin{equation}
  \mathcal{M}_{\mathrm{KKT}} = \mathrm{diag}\bigl(\mathcal{M}_{S},\;
     P_\lambda\bigr), \qquad P_\lambda = \theta\hat M^{-1}
     \ (\text{entries }\theta/\hat m_i),
\end{equation}
with $\mathcal{M}_S$ the problem's block preconditioners rebuilt on the
consistently augmented blocks of $S_\theta$, and $P_\lambda$ the inverse
of the Schur approximation $\hat M/\theta$. (With exact $\mathcal{M}_S =
S_\theta^{-1}$ and exact Schur inverse the preconditioned KKT matrix would
have only the three eigenvalues $1$, $(1\pm\sqrt5)/2$; the approximations
above spread these into clusters.) The scaling is easy to get backwards:
interface mass entries scale like $h^{d-1}$, so using $\theta\hat M$ instead
of $\theta\hat M^{-1}$ mis-scales the block by $\theta^2\hat m^2$ (about
$10^4$ at $h = 0.3$); the correct form cut 79.4k to 22.9k iterations while
moving the endpoint only at solver tolerance --- the two properties any
change to this block should be re-verified against. The projector removes
the slip null pairs (\S\ref{sec:slip-null}) extended by a zero multiplier
component, which the constraint row annihilates.

\paragraph{Algorithm (KKT path).}
\begin{enumerate}
\item Assemble $\mathcal{A}$, $\eps Q$, the kernel $N$ and $\hat M$; form
  $\mathcal{C}$ and the augmented blocks $S_\theta$; build
  $\mathcal{M}_{\mathrm{KKT}}$ and the projector.
\item Solve \eqref{eq:kkt} by projected MINRES; set $\delta\leftarrow U$.
\item For $k = 1,\dots,K$ (refinements, bounded by \code{SetConstraint}'s
  count): by identity (i) of \S\ref{sec:tikhonov} the physical residual is
  $[\eps Q\delta;\,0]$; solve \eqref{eq:kkt} with that right-hand side
  (zero constraint and potential loads) for the increment
  $(\delta_{\mathrm{new}}, \delta_\lambda)$, accumulate $U\leftarrow U +
  \delta_{\mathrm{new}}$, $\lambda\leftarrow\lambda + \delta_\lambda$, and
  set $\delta\leftarrow\delta_{\mathrm{new}}$. The inexact
  schedule of \S\ref{sec:al} may be applied to these refinement solves.
\item Return $U$ and $\lambda$ (\code{KKTMultiplier()}).
\end{enumerate}

\paragraph{Enforcement is weak.} $\mathcal{C}U = 0$ zeroes the jump
against the multiplier space, so the $L^2$ jump of the solution is the
projection remainder --- the same discrete constraint floor a converged AL
iteration reaches; the difference between the paths is the route, not the
destination. The multiplier is physical (the interface normal-traction
perturbation paired with $\bnu$) and cross-checks the AL path's
accumulated $w$.

\paragraph{Multiplier space.} Equal order (the displacement's) is the
strictest choice; one order lower is the \emph{mortar}-style choice (as in
mortar methods \cite{Wohlmuth2001}, where the multiplier space is chosen
coarser than the trace space), improving the inf--sup margin of the trace
pairing and shrinking the multiplier block (measured $\sim1.6\times$
cheaper at an endpoint shift of $4\times10^{-4}$, well below mesh error);
it is the recommended benchmark setting. $\theta = 100$ is the default in
both paths ($\theta = 300$ gains only 5\,\%).

\paragraph{Head-to-head} (fluid core, $h = 0.3$, order 2):
\begin{center}
\small
\begin{tabular}{@{}lrr@{}}
\toprule
solver & iterations & wall \\
\midrule
AL, 3 sweeps, exact & 4.4k & 76 s \\
AL, 8 sweeps, exact & 11.4k & 200 s \\
KKT, order-1 multiplier, 3 refinements, exact & 10.7k & 178 s \\
KKT, same, inexact refinement sweeps & 7.4k & 124 s \\
\bottomrule
\end{tabular}
\end{center}
KKT and AL agree within discretisation error everywhere (at degree 0 to
$2\times10^{-7}$ with identical error against \code{pyslfp}); endpoint
differences between sweep counts and paths (up to $2\times10^{-2}$ at
degree 1) sit below the mesh error, and 3-sweep AL runs are not worse than
8-sweep ones against the external reference. The KKT endpoint is invariant
to $\theta\in[1,300]$ at the $5\times10^{-10}$ level, as the formulation
demands (the AL endpoint moves with $\theta$ through its constraint floor):
a sharp regression check on the constraint blocks.

\subsection{Null space}
\label{sec:slip-null}

Common translations $(\vt,\vt,0)$ and \emph{independent} mapped rotations
of shell and core (a frictionless axisymmetric interface transmits no
torque), plus the 2-D potential constant. Translations of either component
alone are not null (gravity couples them). With the tapered extensions the
translations are near-null rather than exact. The rigid pairs are near-null
under the full operator including both penalties, since $\bnu$ and $\vb$
are radial against tangential rotations. On an elliptical interface the
independent-rotation modes disappear (torque is transmitted through normal
forces alone).

\subsection{Verification}

The interface form \eqref{eq:BSigma} matches the exact energy of
constrained broken motions to $2\times10^{-4}$ (identity and curved radial
$\varphi_e$; diagonal and polarised second differences); the gravity
pieces to $\le3\times10^{-4}$ on families with a rigid, stress-free solid,
including the jointly annihilated axisymmetric relabelling; $G_\Sigma$ to
$3\times10^{-4}$, resolved at $30\times$ the residual. Solution-level, on
the two-layer hydrostatic disc at order 2:
\begin{itemize}
\item single-valued slip versus the welded gauged solution (tangential
  continuity being an admissible gauge for the barotropic fluid):
  $1.1\times10^{-3}$ (solid displacement), $7\times10^{-4}$ (solid-region
  potential); two fluid extensions (taper powers 2, 3) agree to
  $5\times10^{-6}$, the built-in gauge-invariance test;
\item reduction to the Dahlen class under \eqref{eq:zetaphi}, degree-2 load:
  $6\times10^{-3}$ (displacement), $5\times10^{-3}$ (potential) --- two
  formulations sharing no interface machinery;
\item gravity and pressure off: the gravity-free sliding interface, matching
  the condensed cavity reference to $2\times10^{-5}$;
\item broken versus single-valued head-to-head: $7\times10^{-5}$
  (displacement), $5\times10^{-5}$ (potential), both constraints contracting
  by $\sim10^2$.
\end{itemize}
Note that $\zeta^1 = \varphi^1 + \tilde\vv\cdot\grad\zeta^0$ carries the
extension field off the solid, so it is an observable only on the solid
region. The mapped broken solver reproduces the
identity-description solution under an interface-fixing twist relabelling
to $5\times10^{-4}$ and $2\times10^{-4}$. Serial and parallel agree at
$10^{-5}$ on 1/2/4 ranks.

\paragraph{Published treatments.} Checked against this derivation, the
natural-configuration interface terms of AC18 \cite{AC18} (their eq.~121)
are exactly \eqref{eq:BSigma} at $\varphi_e = \mathrm{id}$ with
$\varpi^{(0)}$ identified with $\pi$; the linearised slip of Maitra \&
Al-Attar \cite{MaitraAlAttar2024} agrees in its exact theory, but their
linearised $\mathbf{Q}$-terms (eqs.~B5--B6) miss a factor $J^{(0)}$,
invisible at a natural configuration.

%==========================================================================
\section{The viscoelastic time layer}
\label{sec:visco}
%==========================================================================

The layer models generalised Maxwell (Prony-series) bodies, isotropic or
anisotropic: with strain $\eps$, a relaxed tensor $C_\infty$, and for each
\emph{branch} $k$ a relaxable tensor $C_k$, a relaxation time $\tau_k$ and a
symmetric-tensor internal variable $m_k$,
\begin{equation}
  \sigma = C_U\eps - \sum_k C_km_k, \qquad C_U = C_\infty + \sum_kC_k,
  \qquad \dot m_k = (\eps - m_k)/\tau_k,\ m_k(0) = 0 ,
\end{equation}
$C_U$ being the unrelaxed (instantaneous) tensor. For the isotropic body
$C_k = 2\mu_kP_{\mathrm{dev}}$ ($\mu_k$ the branch shear modulus,
$P_{\mathrm{dev}}$ the projector onto deviatoric tensors) and the internal
variables are trace-free; for an anisotropic body which part of $C$ relaxes
is a modelling choice of the branch coefficient. The elastic operator is
assembled with the \emph{unrelaxed} modulus and the branches appear only as
forces through the problem interface of \S\ref{sec:common}:
\begin{equation}
  K_U\vu = f_{\mathrm{ext}}(t) + \sum_kB^T(C_km_k), \qquad
  \dot m_k = (D\vu - m_k)/\tau_k,
\end{equation}
with $K_U$ the stiffness matrix assembled with $C_U$, $f_{\mathrm{ext}}$
the external load, $\Phi_I$ the tensor-valued basis functions of the
internal-variable space and $\phi_J$ the displacement basis functions,
$B_{IJ} = \int\Phi_I:\eps(\phi_J)$ assembled once with unit coefficient
(material weighting $C_km_k$ applied pointwise at the internal-variable
nodes), and $D$ the strain map from displacement to internal-variable
nodes (below). Self-gravitation, fluid cores and coupled unknowns live inside
the problem, so the layer runs unchanged on coupled problems and never sees
a potential; in the self-gravitating class only the displacement block is
reassembled. \code{ReferentialElasticRheology} has no branches (its
relaxation weights are no-ops), so the referential classes are elastic.

\paragraph{Internal variables.} On an $L^2$ (discontinuous) nodal space of
the displacement mesh; the strain map $D$ is either the Galerkin projection
$(G_E^{-1}\otimes M_{L^2}^{-1})B$ (default; adjoint-consistent), with
$M_{L^2}$ the scalar $L^2$ mass matrix (inverted exactly element by
element) and $G_E$ the Frobenius metric $(G_E)_{cc'} = E_c:E_{c'}$ of the
non-orthonormal basis tensors $E_c$ of the component convention, or nodal
interpolation. The internal order must resolve $\eps(\vu)$ exactly ($p-1$
on simplices, $p$ on tensor-product elements, $p$ the displacement order),
so that the effective-modulus elimination is exact.

\paragraph{Steppers.} With step $dt$, $h_k = dt/\tau_k$,
$b_k = 1 - (1 - e^{-h_k})/h_k$ (the exponential-trapezoid weight of the end
strain), $\gamma$ the SDIRK diagonal coefficient, and the \emph{relaxation
weights} $\beta_k$ defining the effective modulus $C_\infty + \sum_k\beta_kC_k$
of the solve:
\begin{center}
\small
\begin{tabular}{@{}P{3.4cm}P{5.4cm}P{4.6cm}@{}}
\toprule
Scheme & elastic solve & order, stability \\
\midrule
explicit RK & one per stage, $K_U$ & RK order; $dt\lesssim2.8\tau_{\min}$ (RK4) \\
ETD1 & one per step, $K_U$ & 1st; unconditionally stable \\
exponential trapezoid & one per step, weights $\beta_k = 1 - b_k$ & 2nd; exact for strain linear in time \\
backward Euler & weights $\beta_k = 1/(1+h_k)$ & 1st; L-stable \\
SDIRK & as backward Euler, $\gamma dt$ per stage & 2nd--3rd \\
\bottomrule
\end{tabular}
\end{center}
(RK: Runge--Kutta; ETD1: first-order exponential time differencing;
SDIRK: singly diagonally implicit Runge--Kutta; L-stable: stable with
damping of infinitely stiff modes.) Implicit and exponential-trapezoid steps
reassemble the stiffness with the effective modulus (pointwise weights,
since $\tau_k$ varies in space), through \code{SetRelaxationWeights}; the
preconditioner-reuse rule of \S\ref{sec:nograv} keeps the AMG setup across
these reassemblies; Crank--Nicolson is
deliberately absent (not L-stable). With stiff branches SDIRK23 is the best
fixed-step scheme and the adaptive trapezoid the best at tight tolerances.
State-dependent relaxation times $\tau_k = \tau_{k0}F_k(\eps,\sigma,m_k)$
($\tau_{k0}$ a reference time, $F_k$ a pointwise factor)
(e.g.\ a power law, composite diffusion/dislocation creep, after Crawford et
al.\ \cite{Crawford2017} and Simo \& Hughes \cite{SimoHughes}) keep every
global solve linear: a predictor--corrector re-evaluates $\tau$ at the
midpoint state, and adaptive stepping uses the embedded ETD1 companion as a
(conservative) error estimate. \code{CompositeRheology} assigns different
rheologies to different attribute regions of one displacement space, each
branch living on its region's nodes only.

%==========================================================================
\begin{thebibliography}{99}
%==========================================================================

\bibitem{AWahrZhong2013}
G.~A, J.~Wahr \& S.~Zhong, 2013.
Computations of the viscoelastic response of a 3-D compressible Earth to
surface loading: an application to Glacial Isostatic Adjustment in
Antarctica and Canada, \emph{Geophys.\ J.\ Int.}, 192, 557--572.

\bibitem{AC18}
D.~Al-Attar, O.~Crawford, A.~P.~Valentine \& J.~Trampert, 2018.
Hamilton's principle and normal mode coupling in an aspherical planet with
a fluid core, \emph{Geophys.\ J.\ Int.}, 214, 485--507.

\bibitem{AlAttarCrawford2016}
D.~Al-Attar \& O.~Crawford, 2016.
Particle relabelling transformations in elastodynamics,
\emph{Geophys.\ J.\ Int.}, 205, 575--593.

\bibitem{AlAttarTromp2014}
D.~Al-Attar \& J.~Tromp, 2014.
Sensitivity kernels for viscoelastic loading based on adjoint methods,
\emph{Geophys.\ J.\ Int.}, 196, 34--77.

\bibitem{AW10}
D.~Al-Attar \& J.~H.~Woodhouse, 2010.
On the parametrization of equilibrium stress fields in the Earth,
\emph{Geophys.\ J.\ Int.}, 181, 567--576.

\bibitem{Crawford2017}
O.~Crawford, D.~Al-Attar, J.~Tromp \& J.~X.~Mitrovica, 2017.
Forward and inverse modelling of post-seismic deformation,
\emph{Geophys.\ J.\ Int.}, 208, 845--876.

\bibitem{DahlenTromp1998}
F.~A.~Dahlen \& J.~Tromp, 1998.
\emph{Theoretical Global Seismology}, Princeton University Press.

\bibitem{EnglHankeNeubauer}
H.~W.~Engl, M.~Hanke \& A.~Neubauer, 1996.
\emph{Regularization of Inverse Problems}, Kluwer.

\bibitem{FortinGlowinski}
M.~Fortin \& R.~Glowinski, 1983.
\emph{Augmented Lagrangian Methods}, North-Holland.

\bibitem{FriedmanSchutz1978}
J.~L.~Friedman \& B.~F.~Schutz, 1978.
Lagrangian perturbation theory of nonrelativistic fluids,
\emph{Astrophys.\ J.}, 221, 937--957.

\bibitem{GolubGreif2003}
G.~H.~Golub \& C.~Greif, 2003.
On solving block-structured indefinite linear systems,
\emph{SIAM J.\ Sci.\ Comput.}, 24, 2076--2092.

\bibitem{Hestenes1969}
M.~R.~Hestenes, 1969.
Multiplier and gradient methods, \emph{J.\ Optim.\ Theory Appl.}, 4,
303--320.

\bibitem{Huang2023}
P.~Huang, R.~Steffen, H.~Steffen, V.~Klemann, P.~Wu, W.~van der Wal,
Z.~Martinec \& Y.~Tanaka, 2023.
A commercial finite element approach to modelling Glacial Isostatic
Adjustment on spherical self-gravitating compressible earth models,
\emph{Geophys.\ J.\ Int.}, 235, 2231--2256.

\bibitem{Lardy1975}
L.~J.~Lardy, 1975.
A series representation for the generalized inverse of a closed linear
operator, \emph{Atti Accad.\ Naz.\ Lincei Rend.\ Cl.\ Sci.\ Fis.\ Mat.\
Natur.\ (8)}, 58, 152--157.

\bibitem{Latychev2005}
K.~Latychev, J.~X.~Mitrovica, J.~Tromp, M.~E.~Tamisiea, D.~Komatitsch \&
C.~C.~Christara, 2005.
Glacial isostatic adjustment on 3-D Earth models: a finite-volume
formulation, \emph{Geophys.\ J.\ Int.}, 161, 421--444.

\bibitem{MaitraAlAttar2021}
M.~Maitra \& D.~Al-Attar, 2021.
On the stress dependence of the elastic tensor,
\emph{Geophys.\ J.\ Int.}, 225, 378--415.

\bibitem{MaitraAlAttar2024}
M.~Maitra \& D.~Al-Attar, 2024.
On the elastodynamics of rotating planets,
\emph{Geophys.\ J.\ Int.}, 237, 1301--1338.

\bibitem{MarsdenHughes}
J.~E.~Marsden \& T.~J.~R.~Hughes, 1983.
\emph{Mathematical Foundations of Elasticity}, Prentice-Hall.

\bibitem{Martinec2000}
Z.~Martinec, 2000.
Spectral-finite element approach to three-dimensional viscoelastic
relaxation in a spherical earth, \emph{Geophys.\ J.\ Int.}, 142, 117--141.

\bibitem{Powell1969}
M.~J.~D.~Powell, 1969.
A method for nonlinear constraints in minimization problems, in
R.~Fletcher (ed.), \emph{Optimization}, Academic Press.

\bibitem{SimoHughes}
J.~C.~Simo \& T.~J.~R.~Hughes, 1998.
\emph{Computational Inelasticity}, Springer.

\bibitem{Spada2011}
G.~Spada et al., 2011.
A benchmark study for glacial isostatic adjustment codes,
\emph{Geophys.\ J.\ Int.}, 185, 106--132.

\bibitem{TrompMitrovica1999}
J.~Tromp \& J.~X.~Mitrovica, 1999.
Surface loading of a viscoelastic earth---I. General theory,
\emph{Geophys.\ J.\ Int.}, 137, 847--855.

\bibitem{Valette1986}
B.~Valette, 1986.
About the influence of pre-stress upon adiabatic perturbations of the
Earth, \emph{Geophys.\ J.\ R.\ astr.\ Soc.}, 85, 179--208.

\bibitem{Valette1991}
B.~Valette, 1991.
Gravito-elastodynamics of a pre-stressed elastic earth,
\emph{Geophys.\ J.\ Int.}, 104, 555.

\bibitem{Wohlmuth2001}
B.~I.~Wohlmuth, 2001.
\emph{Discretization Methods and Iterative Solvers Based on Domain
Decomposition}, Springer.

\bibitem{WoodhouseDeuss2007}
J.~H.~Woodhouse \& A.~Deuss, 2007.
Theory and observations --- Earth's free oscillations, in
\emph{Treatise on Geophysics}, vol.~1, Elsevier.

\bibitem{WuNi1996}
P.~Wu \& Z.~Ni, 1996.
Some analytical solutions for the viscoelastic gravitational relaxation of
a two-layer non-self-gravitating incompressible spherical earth,
\emph{Geophys.\ J.\ Int.}, 126, 413--436.

\bibitem{WuPeltier1982}
P.~Wu \& W.~R.~Peltier, 1982.
Viscous gravitational relaxation, \emph{Geophys.\ J.\ R.\ astr.\ Soc.}, 70,
435--485.

\bibitem{CitcomSVE2025}
T.~Yuan, S.~Zhong \& G.~A, 2025.
CitcomSVE-3.0: a three-dimensional finite-element software package for
modeling load-induced deformation and glacial isostatic adjustment for an
Earth with a viscoelastic and compressible mantle,
\emph{Geosci.\ Model Dev.}, 18, 1445.

\bibitem{Yu2025}
Z.~Yu, D.~Al-Attar, F.~Syvret \& A.~J.~Lloyd, 2025.
Application of first- and second-order adjoint methods to glacial isostatic
adjustment incorporating rotational feedbacks,
\emph{Geophys.\ J.\ Int.}, 240, 329--348.

\bibitem{Zhong2003}
S.~Zhong, A.~Paulson \& J.~Wahr, 2003.
Three-dimensional finite-element modelling of Earth's viscoelastic
deformation: effects of lateral variations in lithospheric thickness,
\emph{Geophys.\ J.\ Int.}, 155, 679--695.

\end{thebibliography}

%==========================================================================
% NOTES TO THE AUTHOR
% Points where the repository's sources disagree, are stale relative to one
% another, or are unclear. No resolution has been invented in the text above;
% where a choice had to be made, the choice is stated here.
%
% 1. Warm or cold refinement solves. doc/gauged_fluid.md section 4 says each
%    Tikhonov refinement is "one further SolveLinearSystem per refinement,
%    cold-started"; doc/gauge_penalty_iteration.tex (section 3.3, and step 3
%    of its algorithm) says each refinement solve is warm-started from the
%    previous iterate. The text above avoids stating either for the
%    refinement solves themselves.
%
% 3. Mapped single-valued slip. The slip-class doxygen lists the mismatch
%    gravity pieces at phi_e = id as a limitation with "mapped variants
%    deferred"; doc/slip_interface.tex (final section) says that by decision
%    they are not implemented, the single-valued assembly refuses a
%    non-identity map (IsIdentity guard) and the broken organisation is the
%    primary for mapped backgrounds. The text follows slip_interface.tex.
%
% 4. The AL multiplier and pi. doc/gravitating_elasticity.md section 5.3
%    (implementation map) and slip_interface.tex section 3 state that the
%    converged multiplier "must approach pi"; slip_interface.tex's working
%    note (ii) says the converged multiplier of the LINEARISED solve is the
%    first-order traction varpi^(1), the pi diagnostic belonging to a
%    first-variation/nonlinear solve. The text follows the working note.
%
% 5. Covariance of the slip constraint forms. doc/mappings.md says the
%    relabelled identity through the slipping interface is informational
%    (u 1.2e-2) because "the interface constraint forms are not yet
%    certified covariant"; doc/kkt_slip_solver.md says the KKT constraint
%    row is covariant by construction (Nanson-exact flux kernel), and
%    slip_interface.tex reports the mapped broken solver verified to 5e-4 /
%    2e-4 by the tier-(ii) relabelling. It is unclear which forms remain
%    uncertified (the AL normal-normal penalty with its 1/|nu|, the
%    scalar-jump kernels?) and whether the 1.2e-2 refers to the AL path only.
%
% 6. Stale "open" statements. doc/gauged_fluid.md section 5 and
%    doc/gauge_penalty_iteration.tex section 4.5 still list the three-block
%    self-gravitating slip solver and the interface gravity terms as open;
%    slip_interface.tex records them as built and verified. Likewise
%    gravitating_elasticity.md sections 5.2-5.3 carry the provisional
%    C_Sigma coupling and "flux-variation x jump" term that
%    slip_interface.tex shows to vanish / be an artefact. The text follows
%    slip_interface.tex throughout. slip_interface.tex itself is still
%    headed "FOR REVIEW".
%
% 7. Literature. The notes refer to "Dahlen's treatment" without a primary
%    Dahlen reference; only Dahlen & Tromp (1998) and AW10's rendering are
%    cited here. Bibliographic details of Tromp & Mitrovica (1999a), Martinec
%    (2000), Zhong et al. (2003), Wu & Ni (1996), Friedman & Schutz (1978),
%    Golub & Greif (2003), Crawford et al. (2017), Marsden & Hughes and
%    Simo & Hughes are not in the repository and were supplied from general
%    knowledge -- please check. (BenchmarkPapers/137-2-469.pdf
%    is Martinec 1999, GJI 137, 469, which the notes do not cite by year.)
%    Valette (1991)'s title is as given in gravitating_elasticity.md.
%
% 8. RelabelledBackground's tensor law is quoted from background.hpp as
%    J Q^{-T} (C o xi) Q^{-1}; Q_xi is not defined in any note.
%    gravitating_elasticity.md gives the law for the FIRST tensor A only.
%    Relatedly, the mapping layer's convention ("nothing absorbed into
%    relabelled densities", mappings.md section 1) and the background
%    layer's absorbed rho~ = J rho o xi (AC18 convention) are two different
%    roles of a mapping; the text explains the distinction, but no single
%    note states it in one place.
%
% 9. Fluid-row mass term in the single-valued slip. The class doxygen calls
%    it the rho sym(D g0) mass term; slip_interface.tex writes
%    rho grad grad zeta0. Identical for an exact gradient; with a discrete
%    zeta0 the symmetrised gradient of the discrete g0 is presumably what is
%    assembled -- worth confirming.
%
% 10. Gauge penalty operating point (updated). The two-sided certification
%     (kkt_slip_solver.md section 7, and the coordinator's figures: bias
%     2.6e-2 in tidal h2 at eps = 1e-1; h2 = -3.17 at eps = 1e-4) fixes the
%     operating point at eps ~ 1e-2, as the text now states. doc/gauged_fluid.md
%     section 1 still says "eps ~ 1e-2..1e-3", and the two measured values
%     above are not recorded in any repository note found (only the 4.5x per
%     decade and the eps = 1e-4 semi-convergence are, in kkt_slip_solver.md);
%     the note may want updating.
%
% 11. AL sweep default. The class default is 8 sweeps (theta = 100);
%     kkt_slip_solver.md finds 3 sweeps not worse against pyslfp at h = 0.3
%     and says the default oversolves there, to be re-examined at finer h.
%
% 12. Viscoelastic coverage. ReferentialElasticRheology has no branches, so
%     the viscoelastic layer applies to the traction/clamped/mixed classes
%     with Maxwell rheologies; it is not stated which problem class the
%     viscoelastic benchmark family uses.
%
% Added in the revision pass (definition audit and linear-algebra detail):
%
% 14. Definiteness of the mixed block system (resolved). The BlockCG doxygen
%     in mixed_problem.hpp now states that the assembled symmetric system is
%     congruent to a positive-definite one for a stable body, which the text
%     adopts. doc/self_gravitation.md section 3 and the BlockMINRES wording
%     still call it "symmetric and indefinite (a saddle point)" -- stale.
%
% 15. Inexact-sweep schedule. The notes say "relaxed to loose_rel and
%     tightened geometrically to the solver's relative tolerance, the final
%     sweep always at full tolerance"; the exact formula is not documented.
%     The text gives tau_k = tau_loose (tau/tau_loose)^{k/(K-1)} as an
%     example only.
%
% 16. 2-D potential-load compatibility. "Subtracting a uniform flux through
%     the outer boundary" is written as f <- f - (1^T f/|dB|) l with
%     l_i = int_dB phi_i dS. This is my reading of the members ones_,
%     L_outer_, outer_length_; the notes give no formula.
%
% 17. The "1/eps" amplification of an unprojected preconditioner's null
%     component (self_gravitation.md section 4, null_space.md) uses an eps
%     that is not defined in either note; the text quotes it as such.
%
% 18. The KKT preconditioner paragraph adds the standard fact that with
%     exact block inverses the block-diagonally preconditioned saddle-point
%     matrix has eigenvalues 1 and (1 +- sqrt 5)/2 (Murphy, Golub & Wathen,
%     2000). This is general knowledge, not from the repository, and is not
%     cited in the bibliography. Likewise the Woodbury-based contraction
%     analysis of the AL sweeps and the Golub-Greif clustering argument are
%     standard derivations written out here; the notes state only their
%     conclusions (||A_S||/(||A_S|| + theta) and Schur ~ M_lambda/theta).
%
% 19. KKT gauge-refinement right-hand side: the text infers from identity
%     (i) of the Tikhonov recurrence that each KKT refinement solves for
%     [eps Q delta; 0]. kkt_slip_solver.md says only that the outer loop
%     "consists of the fluid-gauge Tikhonov refinements alone"; the exact
%     residual the code forms is not documented.
%
% 20. Contraction-rate tripwire scope. The 0.9 / 0.2 thresholds are found
%     only in the mixed class's GaugeRefine() override (src/mixed_problem.cpp);
%     the base-class and referential-class refinements appear not to carry
%     them. The text attributes the tripwire to the mixed class only.
%
% 21. Gauge-KKT measurements (30k vs ~800 iterations, stationarity 1e-5)
%     are quoted from kkt_slip_solver.md section 7; it reports them for the
%     2-D three-layer lab and the 3-D fluid_core h = 0.3 case jointly, so
%     which variant gave which number is not separable from the note.
%==========================================================================

\end{document}
